% NAL 2339 — lecture modernisée du volume entier.
%
% Pass: Opus 5.5 (claude-opus-5-5), 2026-09-23 — first pass, unchecked against
% the views by a human, and an interpretation rather than a record. Where this
% file and the transcriptions differ in substance, the transcriptions
% (batch-01.fr.tex to batch-03.fr.tex) are the record.
%
% Sources read: the three batch transcriptions of this volume, views 1–60, in
% full, twice — once for the mathematics, once for the apparatus (every
% \uncertain, \ill, \note, \marginal, \struck, \add and \folio) — with their
% header comments and the coordinator's reconciliation block closing each
% header, and the NAL 2339 section of .claude/skills/transcribe/references/
% hand.md. The volume is completely transcribed. Views 61–66 are binding
% (endpaper, back board, spine lettered MÉMOIRES DE FERMAT, edges): batch 4
% has no file because it holds nothing to transcribe, not because it is
% missing. Views 1–8 (binding, flyleaves, the library's title leaf, a blank)
% and 36, 39, 52–60 (blanks, endleaves, endpaper) are recorded in the headers
% and not transcribed. No image was consulted. No printed edition was
% consulted: not Tannery and Henry's Œuvres de Fermat, not Apollonius,
% Pappus, Eutocius or Viète. Where the reading states a fact the leaves do not
% carry (Eutocius's report of Menaechmus, the content of Apollonius III.16,
% Wantzel 1837, the death of Descartes), it says so in a footnote as general
% knowledge, checked against no edition; III.16 is in addition verified
% analytically.
%
% Map (settled by the reconciliation notes; library foliation = red crayon,
% 1–24, v. 9 = f. 1r … v. 51 = f. 22r, ff. 23–24 blank on v. 53–56):
%   v. 9–26   « Ad locos planos & solidos Isagoge », ink pagination 1–18
%   v. 27–28  « Copie d'un escrit enuoyé par le R. Pere Marçenne a des Cartes »,
%             « Methodus ad disquirendam maximam & minimam »
%   v. 29–30  « De tangentibus Linearum curuarum », same hand, continues it
%   v. 31–35  « Appendix ad Isagogen topicam … per locos », second hand, foot
%             numbers 43–45
%   v. 37–38  headed « Mr. fermat » (paler, perhaps later): three-line locus;
%             address « Pour Mons.r Carcaui [doubtful] … »; its figure on v. 40
%   v. 41–48  « [Nouus] Secundarum et Vlterioris ordinis Radicum In Analyticis
%             Vsus A Domino de Fermat ad Dominum De Carcaui Die 20a. Aprilis
%             Anno 1650 Missus », with its « Appendix Ad superiorem Methodum »
%             (v. 44–48)
%   v. 49–51  « Lettre de Mons.r De fermes A Monsieur De Carcaui De Castres Le
%             20e auril 1650. », unsigned, ends « &c. »
%
% Attribution. Fermat is named on the leaves in three headings only (v. 37,
% 41, 49) and on the library's 1888 title leaf (v. 7). The reading says
% exactly that, attributes no hand, and does not extend the name to the
% unnamed pieces (v. 9–35): there it says « l'écrivain », « l'auteur de
% l'écrit », « le texte ». That the whole is Fermat's is the catalogue's title,
% presented as such. The link of the Appendix of v. 31–35 to the Isagoge of
% v. 9–26, of the letter to the Latin piece, and of the letter's « methode de
% Tangentibus » to v. 29–30 are stated as inferences (the last one is not
% drawn).
%
% Dating. \dating is empty, as in the catalogue. The only dates are the
% leaves' own, 20 April 1650, on v. 41 and v. 49; they are reported for those
% two pieces and extended to no other. Nothing is dated from Mersenne's or
% Descartes's names; the letter's « a la memoire d'vn si grand … homme » is
% reported as its wording.
%
% Notation translated, each with a note at first use: A → x, E → y (in the
% Methodus E → e, in the tangent piece A → a, E → e; in the letter A, the
% ordinate, → y and E, the abscissa, → x); consonants kept; dimension marks
% (pl., sol., ppl., q^dto) dropped with the law of homogeneity, explained once;
% « Zq » where the leaf writes it → Z²; « in » → product; Aq, Eq, Aqq, A cubus
% → powers; '' and « bis » → 2; « æquatur », « Æle », « Æ. » and the double
% stroke of v. 11 → =; « adæquatur » → ≈ (explained, then replaced by a
% rigorous statement); Lat. Cub., Lat. quad. quad., Lat. quad. Cub. → cube,
% fourth, fifth roots; proportions → quotients; rectangulum XYZ → XY·YZ;
% applicata / ordinata → ordonnée; rectum latus → paramètre; donnée d'espèce →
% de forme donnée; applicatio → division.
%
% Modern names attached, each where the statement coincides: coordinates
% (oblique / Cartesian); line, conic, degenerate conic; hyperbola referred to
% its asymptotes; completing the square; discriminant b² − 4ac, only for the
% leaf's cases (and via the shear that the « triangle given in species » is);
% derivative, stationary point, Fermat's theorem on interior extrema (named as
% a later name); subtangent; solution of cubics and quartics by intersecting
% conics; depressed quartic; resolvent cubic (for Viète's « parapleroses »);
% two mean proportionals, duplication of the cube; Menaechmus (footnote only,
% general knowledge); Wantzel 1837 (footnote, the leaf proves nothing);
% three-line locus / Pappus's problem; bitangent pencil of conics; pole and
% conjugate diameter; Apollonius III.16 as the leaf cites it, verified by the
% power of a point with respect to a conic; elimination, resultant (up to an
% extraneous factor; Sylvester not named, the procedure is not his);
% rationalisation of radicals; overdetermined systems and gcd; circular
% sections (cyclic planes) of a quadric cone; n-ellipse and its normal.
%
% Departures from the leaf, each footnoted where it occurs: v. 10 the condition
% « neutra … quadratum prætergrediatur » restated as total degree ≤ 2; v. 11
% « EI » for ZI (the margin's « forsan ZI » followed); v. 12 the proposition on
% lines restated with its sign and degenerate cases; v. 13–14 the hyperbola
% (x−S)(R−y) = D − RS with D = RS and D < RS added; v. 14 « uD siue ZI æquetur
% R » (a vertical side of the rectangle, ND, ou or Zp, is needed; the margin's
% « forsan Zp » followed); v. 14–15 homogeneous case restated as a pair of
% lines through N; v. 16 « applicatæ … parallelæ NP » (NZ needed) and « sub D
% in IP » (NP needed), « Parobolen » spelling; v. 19 circle with a²/4 + b²/4 − c
% ≤ 0 added; v. 20–22 degenerate cases of ellipse and hyperbola added; v. 22
% « Punctum igitur Z » for I; v. 24 the corrector's NZ, NR followed; v. 25
% « species » read as the squares, the proposition sharpened to an ellipse;
% v. 25–26 the example is not an instance of the stated proposition, the
% ratio must exceed 4, and the mirror circle is added; v. 27–30 adæqualitas
% replaced by an exact statement, « Nec vnquam fallit methodus » qualified
% (stationary non-extrema, singular points); v. 29 the inequality needed on
% both sides of B; v. 31–35 missing signs supplied (v. 31, 32, 33), « Bq in Aq »
% for « Bq in Aq bis » (v. 34), « æquabitur Z Sol. in A + d ppl. » (v. 33) read
% as D − Zx, « VZ » for nZ (v. 33), extraneous intersections at x = 0 (v. 33,
% 35), « circulum creari Semper » (v. 35) qualified, v. 34 « Z Sol. in d » kept
% as a constant; v. 37 « fertur » read as « secetur », « foe » for foc, the
% analysis replaced by a direct synthesis; v. 42 N²E² for N²E³; v. 43 the final
% substitution carries the extraneous factor (N² − Bx)²; v. 44 « ZA² » read as
% ZA² − A³ (the leaf's second writing); v. 46 the claim that no approximation
% is possible while radicals remain is qualified; v. 47 the abundant example
% needs a single common root.
%
% Steps supplied (the reading's own computations, checked symbolically): the
% expansions of v. 13, 19, 22–24; the shear and the classification table;
% the circle of v. 26 and its bound 4; the rigorous forms of v. 28 and v. 30;
% the general cubic of v. 31; VM as the second mean (v. 33); the conic type
% by the position of V (v. 37); the resultant P² + D³P + 3DPQ − Q³ of degree 6
% (v. 43) and the one-step Euclidean alternative y = (P + DQ)/(Q + D²); the
% sextic (D³ − Zx² − B²x)³ = 27D³(x³ + B²x)(Zx² − x³) completing v. 45; the
% linear root of the abundant example (v. 47); the rectangle's compatibility
% condition (v. 47); the cyclic planes (v. 48); the normal to the 4-ellipse
% (v. 51).
%
% Left unread or unreconstructed: « m » raised (v. 21); the struck « quando »
% (v. 18); « Symetria climatismus Vietæa », « huic », « sese » (v. 44);
% « Emanuerit », « exulerre » (v. 45); « Exposuit » (v. 47); the two words
% « Oti di » (v. 48); the doubtful term « − D²BA » of the letter's curve and
% « B pp^l » (v. 49); the six-point construction of the cone (v. 48) and the
% extensions to curves on surfaces, which the leaf states without computation;
% the letter's tangent itself, whose elimination the leaf does not carry out;
% the pencil note « Mr des Cartes f. 347 » (v. 30); the content of Apollonius's
% plane loci cited on v. 12, 19–20, 25.
\documentclass[11pt,a4paper]{article}
% The preamble shared by every transcription of the mathematicians' manuscripts in Gallica.
%
% It rests on one idea: **uncertainty compounds, it does not resolve.** A
% transcription that smooths over a doubtful word has destroyed the only thing
% it was made for — a reader must be able to tell, at every point, what was on
% the page and what was supplied. The apparatus macros below are therefore the
% critical apparatus, not decoration.
%
% The same commands are recognised by `scripts/render.mjs`, which produces the
% site's reading view, and by `scripts/tei.mjs`. A macro added here must be
% added there too, or rendering fails loudly — which is the wanted behaviour.
%
% Comments here are English, like the rest of the repository. The documents
% this preamble sets are French, and that is deliberate: see the skills.

\ProvidesPackage{germain}[2026/09/15 Transcriptions of mathematicians' manuscripts in Gallica]

\usepackage{amsmath,amssymb,amsthm}
\usepackage{mathrsfs}
% Kept from the parent preamble so the renderer's diagram support stays
% available; these manuscripts draw no commutative diagrams, and nothing here uses it.
\usepackage{tikz-cd}
\usepackage[english,french]{babel}
\usepackage{fontspec}

% --- Greek ------------------------------------------------------------------
% The text face stays fontspec's default, Latin Modern, which has no Greek:
% XeTeX drops every glyph it lacks with a line in the log and nothing on the
% page, and the Greek words the manuscripts quote vanished from the PDFs. So
% the Greek and Greek Extended blocks get a XeTeX character class of their own,
% and entering or leaving a run of it switches the family to CMU Serif — the
% same Computer Modern design, with polytonic Greek — and back. Only the
% family changes: a Greek word in \textit or \textbf keeps its shape and series.
% The transitions to and from every other class carry the tokens that class
% has towards an ordinary letter, so French spacing before « ; » or inside
% « » still holds after a Greek word. No grouping, so no brace or \endgroup
% can come out unbalanced; maths is left alone, since XeTeX inserts nothing
% there. Kept identical in `exercices.sty`.
\makeatletter
\newfontfamily\germainGreekfont{cmun}[
  Extension=.otf, NFSSFamily=germaingreek,
  UprightFont=*rm, ItalicFont=*ti, SlantedFont=*sl,
  BoldFont=*bx, BoldItalicFont=*bi, BoldSlantedFont=*bl]
\newXeTeXintercharclass\germain@greek
\def\germain@greekrange#1#2{%
  \@tempcnta=#1\relax
  \loop
    \XeTeXcharclass\@tempcnta=\germain@greek
    \ifnum\@tempcnta<#2\advance\@tempcnta\@ne\repeat}
\germain@greekrange{"0370}{"03FF}
\germain@greekrange{"1F00}{"1FFF}
\def\germain@greekprev{\rmdefault}
\def\germain@greekin{%
  \edef\germain@greekprev{\f@family}\fontfamily{germaingreek}\selectfont}
\def\germain@greekout{\fontfamily{\germain@greekprev}\selectfont}
\def\germain@greekjoin{%
  \XeTeXinterchartokenstate=\@ne
  \@tempcnta=\z@
  \loop
    \ifnum\@tempcnta=\germain@greek\else
      \XeTeXinterchartoks\germain@greek\@tempcnta=\expandafter{%
        \expandafter\germain@greekout\the\XeTeXinterchartoks\z@\@tempcnta}%
      \XeTeXinterchartoks\@tempcnta\germain@greek=\expandafter{%
        \the\XeTeXinterchartoks\@tempcnta\z@\germain@greekin}%
    \fi
    \ifnum\@tempcnta<4095\advance\@tempcnta\@ne\repeat}
% babel-french sets its own classes, and saves and restores the interchar
% state, each time a language is selected: join again after it does.
\addto\extrasfrench{\germain@greekjoin}
\addto\extrasenglish{\germain@greekjoin}
\AtBeginDocument{\germain@greekjoin}
\makeatother

\usepackage{geometry}
\geometry{a4paper,margin=25mm,marginparwidth=28mm}
\usepackage{marginnote}
\usepackage{xcolor}
\usepackage[normalem]{ulem}
\setlength{\emergencystretch}{3em}
\usepackage{hyperref}
\hypersetup{colorlinks=true,linkcolor=black,urlcolor=black,pdfborder={0 0 0}}

% --- Batch metadata ---------------------------------------------------------
% Taken as they stand by the rendering script, which shows them at the head of
% the reading view. They fix what the file claims to cover. The macro names are
% the parent project's, so that the scripts stay shared; \folder here means the
% volume's slug — `fr-9115` — and \pages counts Gallica views.

\newcommand{\folder}[1]{\gdef\grothFolder{#1}}
\newcommand{\batch}[1]{\gdef\grothBatch{#1}}
\newcommand{\pages}[2]{\gdef\grothFirst{#1}\gdef\grothLast{#2}}
\newcommand{\foldertitle}[1]{\gdef\grothFoldertitle{#1}}

% The holder's own shelfmark, printed as they write it: « Français 9115 ».
\newcommand{\shelfmark}[1]{\gdef\grothShelfmark{#1}}

% The dating, copied verbatim from the holder's catalogue — never inferred
% here. « XIXe siècle » is the BnF's; a pass that dates a leaf from its content
% says so in a \note{}, as its own inference.
\newcommand{\dating}[1]{\gdef\grothDating{#1}}

% The status notice, printed in the title block and repeated as a diagonal
% watermark on every page of the PDF. The manuscripts are in the public
% domain; what this line declares is not a right but a quality — an
% unverified first pass by a machine. The skills write the line; a file
% without it gets no watermark, so leaving it out is visible in the output.
\newcommand{\watermark}[1]{\gdef\grothWatermark{#1}}

\gdef\grothFolder{?}\gdef\grothBatch{}\gdef\grothFirst{?}\gdef\grothLast{?}%
\gdef\grothFoldertitle{}\gdef\grothShelfmark{}\gdef\grothDating{}\gdef\grothWatermark{}

% --- The résumé -------------------------------------------------------------
\newenvironment{resume}
  {\small\setlength{\parskip}{0.45em}}
  {\par\medskip\hrule\medskip}

% --- Keywords ---------------------------------------------------------------
% In English, closing the modernised reading's résumé: the modern vocabulary
% under which to search for what the leaves construct. The manifest extracts
% them to tag the volume — this line is the single source of the tags.
\newcommand{\keywords}[1]{\par\smallskip\noindent
  {\footnotesize\textsc{Keywords} --- \textit{#1}}\par}

% --- The critical apparatus -------------------------------------------------

% The Gallica view number, in the margin. This is the anchor that lets a
% disagreement over a formula be settled by looking at *one* image rather than
% a volume of 750. The site uses it to turn the facsimile as the transcription
% is scrolled. A view is one scanned image: usually one side of a leaf.
\newcommand{\page}[1]{%
  \marginnote{\footnotesize\color{gray!70}\textsf{v.\,#1}}%
  \label{page:#1}\ignorespaces}

% The BnF's pencilled foliation, where the leaf carries it and a pass can read
% it: « 198r ». This is what the literature cites — « ff. 198r–208v » — and
% the one line that turns a citation into a view. Recorded, never inferred:
% a folio a pass cannot read is not written.
\newcommand{\folio}[1]{%
  \marginnote{\footnotesize\color{gray!50}\textsf{f.\,#1}}\ignorespaces}

% The modernised reading's coarser counterpart to \page{}: which views a
% section is drawn from.
\newcommand{\pagerange}[2]{%
  \marginnote{\footnotesize\color{gray!70}\textsf{v.\,#1--#2}}%
  \ignorespaces}

% Illegible: never guessed.
\newcommand{\ill}{\textcolor{red!60!black}{[\,\ldots\,]}}

% A doubtful reading: the word is given, and flagged as uncertain.
\newcommand{\uncertain}[1]{\uline{#1}}

% An editorial addition — a missing letter, an implied word.
\newcommand{\add}[1]{\textcolor{blue!50!black}{[#1]}}

% Struck out by the writer. What was crossed out is part of the document.
\newcommand{\struck}[1]{\sout{#1}}

% The transcriber's note, on the state of the page or the mathematics.
\newcommand{\note}[1]{\par\smallskip{\small\color{gray!60!black}\textit{[#1]}}\par\smallskip}

% A marginal note of hers, distinct from the body of the text.
\newcommand{\marginal}[1]{\par\smallskip\noindent\hspace{1em}%
  {\small\color{gray!60!black}#1}\par\smallskip}

% --- Title ------------------------------------------------------------------
% The title block is French, like the documents themselves.

\AddToHook{shipout/background}{%
  \ifx\grothWatermark\empty\else
    \begin{tikzpicture}[remember picture, overlay]
      \node[rotate=52, scale=3.4, align=center, text=gray!18,
            font=\sffamily\bfseries]
        at (current page.center) {\grothWatermark};
    \end{tikzpicture}%
  \fi}

\AtBeginDocument{%
  \begin{center}
    {\large\bfseries Manuscrits de mathématiciens (Gallica) ---
      \ifx\grothShelfmark\empty\grothFolder\else\grothShelfmark\fi}\\[2pt]
    {\itshape\grothFoldertitle}\\[4pt]
    {\small vues \grothFirst--\grothLast
      \ifx\grothBatch\empty\else\ (lot \grothBatch)\fi}\\[2pt]
    \ifx\grothDating\empty\else
      {\small\color{gray!60!black}Datation du catalogue : \grothDating}\\[2pt]
    \fi
    \ifx\grothWatermark\empty\else
      {\small\bfseries\color{red!45!gray}\grothWatermark}\\[2pt]
    \fi
    {\footnotesize\color{gray!60!black}Transcription établie d'après les images IIIF de Gallica.
     Source gallica.bnf.fr / Bibliothèque nationale de France.\\
     Les lectures incertaines sont soulignées, les illisibles notées [\,\ldots\,],
     les ajouts entre crochets ; \textsf{v.} numérote les vues de Gallica,
     \textsf{f.} la foliotation au crayon de la BnF.}
  \end{center}
  \vspace{1em}}

\endinput


\folder{nal-2339}
\shelfmark{NAL 2339}
\pages{1}{66}
\dating{}
\watermark{Lecture automatique\\interprétation non vérifiée}
\foldertitle{Fermat. Mémoires de mathématiques. NAL 2339 — lecture modernisée du volume entier}

\begin{document}

\section*{Résumé}

\begin{resume}
Ce volume de vingt-deux feuillets écrits réunit six pièces latines et une
lettre française de géométrie et d'algèbre, sous le titre que la bibliothèque
lui a donné en 1888 : « Mémoires de Fermat ». Une introduction aux lieux plans
et solides, « Ad locos planos \& solidos Isagoge » (vues 9 à 26) ; une méthode
pour trouver le maximum et le minimum, puis une méthode des tangentes, sous le
titre français « Copie d'un escrit enuoyé par le R. Pere Marçenne a des Cartes »
(vues 27 à 30) ; un « Appendix ad Isagogen topicam » qui résout les équations du
troisième et du quatrième degré par des courbes (vues 31 à 35) ; une
démonstration sur le lieu à trois droites (vues 37 à 40) ; enfin une méthode
d'élimination des « racines secondes », envoyée à Carcavi avec une lettre
d'accompagnement en français, datées toutes deux sur les feuillets du 20 avril
1650 (vues 41 à 51). Le nom de Fermat n'est écrit sur les feuillets que dans
trois titres, aux vues 37, 41 et 49 ; les quatre premières pièces ne portent
aucun nom, et c'est le titre du catalogue qui fait de l'ensemble les mémoires
de Fermat. Aucune main n'est identifiée. Le volume est entièrement transcrit :
après la vue 51, il n'y a que des feuillets blancs et la reliure.

Le problème commun est de rendre générales et presque mécaniques des questions
que la géométrie grecque traitait cas par cas : sur quelle courbe se trouve un
point défini par une condition (un « lieu ») ; où une grandeur atteint son
maximum ; comment mener la tangente à une courbe ; comment construire la racine
d'une équation du troisième ou du quatrième degré ; comment débarrasser une
équation de ses radicaux. Les feuillets disent eux-mêmes ce qu'ils cherchent.
L'Isagoge s'ouvre sur un reproche aux Anciens, qui ont beaucoup écrit sur les
lieux mais « ne les ont pas exprimés assez généralement » : « nous soumettons
donc cette science à une analyse propre, pour ouvrir désormais aux lieux une
voie générale ». Et la lettre commence ainsi : « Je n'ay point voulu differer a
vous Enuoyer ma methode generale pour le desbrouilement des Asymmetries »,
c'est-à-dire des radicaux, en demandant pour seul intérêt « la satisfaction …
d'auoir deuoilé vne matiere qui n'estoit pas cognüe ».

L'idée de l'Isagoge est celle des coordonnées : deux longueurs inconnues,
portées sous un angle donné à partir d'un point fixe, font d'une équation entre
elles une courbe ; si l'équation est du premier degré, c'est une droite ; si
elle est du second, une conique, que l'on reconnaît en complétant les carrés.
L'idée de la méthode des maxima est de comparer la grandeur en $x$ et en $x+e$,
de diviser par $e$, puis d'effacer ce qui contient encore $e$ : on obtient ce
qu'on écrit aujourd'hui $f'(x)=0$. La tangente à la parabole se trouve de la
même façon. L'Appendix coupe une équation du troisième ou du quatrième degré en
deux équations à deux inconnues, deux coniques, dont l'intersection donne la
racine : parabole et hyperbole, deux paraboles, parabole et cercle ; il en tire
les deux moyennes proportionnelles. La pièce envoyée à Carcavi combine deux
équations en deux inconnues pour faire baisser le degré de la seconde jusqu'à ce
qu'une simple division la donne, puis la remplace dans l'une des équations : on
a éliminé une inconnue ; en appelant chaque radical une nouvelle inconnue, on
fait disparaître les radicaux. Les calculs des feuillets ont tous été refaits :
ils sont justes, à quelques signes et lettres près, qui sont signalés ; là où
une démarche n'est pas une démonstration — l'« adéquation », l'analyse du lieu à
trois droites —, la lecture le dit et donne la démonstration.

Ce que ces feuillets construisent s'appelle aujourd'hui la géométrie analytique
et la classification des coniques (par le discriminant $b^2-4ac$) ; la dérivée,
les points stationnaires et le théorème de Fermat sur les extremums intérieurs ;
la sous-tangente ; la résolution des équations par intersection de coniques, les
deux moyennes proportionnelles et la duplication du cube ; le lieu à trois
droites, ou problème de Pappus, avec la proposition III.16 des Coniques
d'Apollonius ; l'élimination et le résultant ; les systèmes surdéterminés ; les
sections circulaires d'un cône.

\keywords{plane and solid loci, analytic geometry, conic sections, classification of conics, adequality, maxima and minima, interior extremum theorem, method of tangents, subtangent, geometric solution of equations, intersection of conics, resolvent cubic, two mean proportionals, duplication of the cube, three-line locus, Pappus's problem, Apollonius's Conics, elimination theory, resultant, elimination of radicals, overdetermined systems, circular sections of a cone}
\end{resume}

\section*{Ce que contient le volume}

\pagerange{1}{66}
Soixante-six vues, une page par vue. Les vues 1 à 6 sont la reliure et les
gardes (une étiquette « LAT NOUV. ACQ. 2339 » sur la contre-garde). La vue 7 est
le feuillet de titre de la bibliothèque, calligraphié : « Mémoires de Fermat »,
le certificat « Volume de 24 Feuillets / Les Feuillets 23 \& 24 sont blancs /
8 Octobre 1888 », « Libri 1848. » et la cote ; la transcription le relève sans
le transcrire, puisqu'il n'est pas du manuscrit. La vue 8 est blanche. Le texte
occupe les vues 9 à 51, avec deux pages blanches (vues 36 et 39) ; suivent la
page blanche de la vue 52, les deux feuillets blancs 23 et 24 (vues 53 à 56),
des gardes (vues 57 à 60) et, aux vues 61 à 66, la reliure : garde, plat, dos
titré « MÉMOIRES DE FERMAT », tranches. C'est pourquoi le quatrième lot de vingt
vues n'a pas de fichier : il ne contient rien à transcrire, et le volume est
transcrit en entier.\footnote{Les trois transcriptions (vues 1–20, 21–40,
41–60) le disent chacune dans la note de réconciliation qui clôt leur en-tête :
« v. 61–66 (batch 4) are binding, so batch 4 has no file ».}

\begin{itemize}
  \item \textbf{Vues 9 à 26 (ff. 1r–9v) — « Ad locos planos \& solidos
    Isagoge »}, « introduction aux lieux plans et solides », en latin, d'une
    cursive posée, avec des réclames au pied des pages et sa propre pagination
    à l'encre, 1 à 18. Aucun nom.
  \item \textbf{Vues 27 et 28 (f. 10r–v) — « Methodus ad disquirendam maximam
    \& minimam »}, sous le titre français « Copie d'un escrit enuoyé par le R.
    Pere Marçenne a des Cartes ».
  \item \textbf{Vues 29 et 30 (f. 11r–v) — « De tangentibus linearum
    curvarum »}, de la même main, qui renvoie à la méthode précédente (« Ad
    superiorem methodum »).
  \item \textbf{Vues 31 à 35 (ff. 12r–14r) — « Appendix ad Isagogen topicam
    continens solutionem problematum solidorum per locos »}, d'une seconde
    main, plus calligraphique, avec au pied des rectos les nombres 43 à 45.
  \item \textbf{Vues 37 et 38 (f. 15r–v), et la figure de la vue 40 (f. 16v)
    — le lieu à trois droites}, sous l'inscription « Mr. fermat », avec au
    dos une adresse « Pour Mons.r Carcaui », de lecture douteuse.
  \item \textbf{Vues 41 à 48 (ff. 17r–20v) — « Novus secundarum et ulterioris
    ordinis radicum in analyticis usus »}, « envoyé par M. de Fermat à M. de
    Carcavi le 20 avril 1650 » selon son titre, suivi d'un « Appendix ad
    superiorem methodum » (vues 44 à 48).
  \item \textbf{Vues 49 à 51 (ff. 21r–22r) — une lettre française}, « Lettre
    de Mons.r De fermes A Monsieur De Carcaui De Castres Le 20e auril 1650 »,
    sans signature.
\end{itemize}

\subsection*{Le nom de Fermat, et les mains}

Le nom de Fermat est écrit sur les feuillets du manuscrit en trois endroits, et
en trois seulement : l'inscription « Mr. fermat » en tête de la vue 37, d'une
encre plus pâle et d'un trait plus fin que le texte, peut-être d'une autre main
et peut-être plus tardive ; le titre latin de la vue 41, « A Domino de Fermat
ad Dominum De Carcaui » ; le titre de la lettre de la vue 49, qui écrit « De
fermes ».\footnote{« De fermes » est lu tel quel, en deux mots, la fin nette ;
la note de réconciliation compte ce titre parmi les trois qui nomment Fermat.
La lettre porte le même destinataire et la même date que la pièce latine de la
vue 41.} Il l'est aussi sur le feuillet de titre de 1888 (vue 7), qui est de la
bibliothèque. Les pièces des vues 9 à 35 — l'Isagoge, la méthode des maxima et
minima, la méthode des tangentes, l'Appendix — ne portent aucun nom d'auteur ;
cette lecture ne leur étend pas celui des titres voisins et parle, pour elles,
de « l'écrivain » ou de « l'auteur de l'écrit ». Que l'ensemble soit de Fermat,
c'est le titre du catalogue, « Fermat. Mémoires de mathématiques. », et c'est
comme tel que cette lecture le présente.

Aucune main n'est attribuée. Les transcriptions décrivent une cursive posée
pour l'Isagoge ; la même main pour les deux pièces des vues 27 à 30 ; « une
seconde main, plus calligraphique », pour l'Appendix, qui est peut-être aussi
celle des vues 37–38, sans que ce soit tranché ; une cursive posée à grandes
initiales pour les vues 41 à 51. Plusieurs indices montrent des copies : le
titre de la vue 27 le dit (« Copie d'un escrit »), une note marginale française
de la vue 29 parle de l'auteur à la troisième personne (« Il a icy nommé B ce
qu'il nomme d par apres »), deux notes « forsan … » de l'Isagoge proposent
d'autres lettres, et, dans l'Appendix, les Z sont d'une encre plus noire,
« comme remplis après coup dans des blancs ».\footnote{Ces trois dernières
observations sont des transcriptions ; en tirer que les pièces sont des copies
est une inférence, que les transcriptions font elles-mêmes avec la même
prudence.} La lettre de la vue 49, elle, dit envoyer « mon original » ; si
les feuillets d'ici sont cet original ou une copie, rien ne le dit.

\subsection*{Les numéros des feuillets, et les dates}

La foliotation de la bibliothèque est une série au crayon rouge, en haut à
droite de chaque recto, blancs compris, de 1 (vue 9) à 24 (vue 55) ; elle finit
sur les deux feuillets blancs 23 et 24, comme le dit le certificat de 1888. Les
titres de section de cette lecture donnent ces feuillets ; le texte cite les
vues. Les autres séries n'en sont pas : une série au crayon gris, décalée d'une
unité aux vues 9 à 19 et irrégulière ensuite ; la pagination à l'encre de
l'Isagoge, 1 à 18 ; les nombres 43, 44, 45 au pied des rectos de l'Appendix,
qui paraissent le compte d'un autre assemblage.\footnote{Les autres marques
sont la cote « Libri 1848 », « R. c. 8070 (88) » et le cachet « ACQ. 8070
(LIBRI) » (vues 9 et 27), les cachets ronds de la Bibliothèque nationale, et,
au crayon gris le long de la marge de la vue 30, d'une autre main, « Mr des
Cartes f. 347 », renvoi dont l'objet n'est pas dit sur le feuillet et que cette
lecture ne cherche pas à identifier.}

Le catalogue ne donne aucune date, et la lecture n'en infère aucune. Les
seules dates sont sur les feuillets : « Die 20a. Aprilis Anno 1650 » dans le
titre de la vue 41 et « Le 20e auril 1650 » dans celui de la vue 49. Elles sont
les dates de ces deux pièces-là, telles qu'elles les portent ; elles ne disent
rien des autres.

\subsection*{La notation des feuillets, et sa traduction}

Toutes les pièces écrivent dans l'algèbre « spécieuse » de Viète : les voyelles
$A$ et $E$ désignent les inconnues, les consonnes $B$, $D$, $G$, $N$, $R$, $S$,
$Z$ les grandeurs données, « in » le produit, « Aq », « Eq », « A cubus »,
« Aqq » le carré, le cube, la quatrième puissance, et deux petits traits après
un terme, ou « bis », le double. L'égalité est dite en mots (« æquatur »,
« æquale », abrégé « Æle » ou « Æ. » aux vues 41 à 48), une fois par deux traits
verticaux (vue 11). Chaque grandeur porte sa dimension : « Z planum » est une
aire, « B solidum » un volume, « D plano-planum » une grandeur de dimension
quatre, et une équation n'ajoute que des termes de même dimension (la loi des
homogènes). Cette lecture traduit : $A \to x$, $E \to y$ (sauf avis contraire,
signalé à chaque fois) ; les consonnes sont gardées ; les marques de dimension
sont abandonnées, une grandeur étant désormais un nombre positif, et « Zq »,
quand le feuillet l'écrit, devient $Z^2$ ; « Lat. Cub. », « Lat. quad. quad. »,
« Lat. quad. Cub. » deviennent $\sqrt[3]{\ }$, $\sqrt[4]{\ }$ et
$\sqrt[5]{\ }$\footnote{« Latus cubicum », côté du cube : la racine cubique ;
« latus quadrato-quadratum », la racine quatrième ; « latus quadrato-cubicum »,
côté du carré-cube, c'est-à-dire de la cinquième puissance : la racine
cinquième.} ; une proportion « $a$ ad $b$ ita $c$ ad $d$ » devient
$a/b=c/d$ ; « rectangulum $AEC$ » devient le produit $AE\cdot EC$.

Dans l'Isagoge, un point fixe $N$ est pris sur une droite fixe $NZM$ ; on porte
$NZ = x$ sur cette droite et, sous un angle donné $\widehat{NZI}$, $ZI = y$. Le
point $I$ a donc pour coordonnées $(x,y)$, cartésiennes si l'angle est droit,
obliques sinon. Le vocabulaire des coniques est celui d'Apollonius : un
« diamètre » et ses « appliquées » ou « ordonnées » (les cordes qu'il coupe en
leur milieu, et leurs moitiés) ; le « côté droit » (« rectum latus »), le
paramètre ; « donné de position », fixe ; un triangle « donné d'espèce », de
forme donnée, c'est-à-dire semblable à un triangle fixe. Enfin les feuillets
ne connaissent que des grandeurs positives ; cette lecture dit, à chaque fois,
quand elle passe aux coordonnées signées.

\section*{« Ad locos planos \& solidos Isagoge » : le programme, la droite, l'hyperbole (ff. 1r–4r, vues 9 à 15)}

\pagerange{9}{15}
L'Isagoge commence par l'histoire du sujet. Pappus, au début de son livre VII,
atteste qu'Apollonius a écrit sur les lieux plans et Aristée sur les lieux
solides ; « mais, ou nous nous trompons, ou la recherche des lieux ne leur
était pas assez aisée : nous le conjecturons de ce qu'ils n'ont pas exprimé
assez généralement un très grand nombre de lieux ». D'où le programme :
« Scientiam igitur hanc propriæ et peculiari analysi subiicimus, vt deinceps ad
locos generalis via pateat. »\footnote{« Nous soumettons donc cette science à
une analyse propre et particulière, pour ouvrir désormais aux lieux une voie
générale. » Les ouvrages cités (le livre VII de Pappus, les Lieux plans
d'Apollonius, les Lieux solides d'Aristée) ne sont pas sur les feuillets, et
cette lecture n'en dit rien de plus que le texte.}

La définition suit, et c'est tout le principe : « chaque fois que, dans
l'équation finale, se trouvent deux quantités inconnues, il y a un lieu, et
l'extrémité de l'une d'elles décrit une ligne, droite ou courbe ». Le lieu est
« plan » si cette ligne est une droite ou un cercle, « solide » si c'est une
parabole, une hyperbole ou une ellipse, « linéaire » pour toute autre courbe.
Aujourd'hui : une équation $F(x,y)=0$ entre les coordonnées d'un point définit
une courbe. L'écrivain annonce que le lieu est plan ou solide « pourvu
qu'aucune des inconnues ne dépasse le carré ».\footnote{« Modo neutra
quantitatum ignotarum, quadratum prætergrediatur » (vue 10). Pris à la lettre,
ce critère porte sur le degré de chaque inconnue séparément, et admettrait un
terme comme $x^2y$, qui donne une cubique. La condition exacte est que
l'équation soit de degré total au plus 2, c'est-à-dire ne contienne que des
termes en $x^2$, $xy$, $y^2$, $x$, $y$ et une constante ; tous les cas que
traite le texte la remplissent, le rectangle $xy$ comptant comme un terme du
« degré du carré ».} L'énoncé moderne est celui-ci : une équation du premier
degré en $(x,y)$ définit une droite ; une équation du second degré définit une
conique — ellipse (ou cercle), parabole, hyperbole — ou bien une conique dite
dégénérée : deux droites, une droite, un point, ou rien. L'Isagoge traite
toutes les coniques non dégénérées, et les droites ; les cas dégénérés, elle ne
les nomme pas, et cette lecture les ajoute là où ils se présentent.

\subsection*{Les équations du premier degré : des droites}

Premier cas, $Dx = By$. Le point $I$ est sur une droite passant par $N$ :
$y/x = D/B$ est un rapport donné, donc le triangle $NZI$, dont l'angle en $Z$
est donné, est de forme donnée, et la droite $NI$ est fixe.

Deuxième cas, $Z - Dx = By$, où $Z$ est une aire.\footnote{« Z pl. $-$ D in A
æqu. B in E », avec, pour signe d'égalité, deux traits verticaux entre lesquels
est écrit « æqu. ».} Posons $R = Z/D$, de sorte que $DR = Z$ ; alors
$y/(R-x) = D/B$. Sur $NZM$, soit $NM = R$ : le point $M$ est fixe, $MZ = R - x$,
et le rapport $ZI/MZ$ est donné ; le triangle $IZM$ est de forme donnée et $I$
est sur une droite fixe passant par $M$.\footnote{Le feuillet écrit « dabitur
ergo ratio MZ ad EI », $E$ étant le nom de la quantité $ZI$ ; une note de la
même encre, dans la marge, propose « forsan ZI ». C'est $ZI$ qu'il faut, et
cette lecture suit la marge.} Et de même, conclut le texte, « dans toute
équation dont certains termes sont affectés de $A$ ou de $E$ » : toute équation
$ax + by = c$, avec $a$, $b$ non tous deux nuls, définit une droite. C'est vrai
dans des coordonnées obliques comme dans des coordonnées rectangulaires.
L'écrivain en tire que tous les lieux à la droite, par exemple la septième
proposition du premier livre des Lieux plans d'Apollonius, s'énoncent et se
construisent désormais plus généralement.

Il ajoute une proposition « très belle, que nous avons découverte par ce
moyen » : si, d'un point, on mène à des droites données de position, en nombre
quelconque, des droites sous des angles donnés, et que la somme des rectangles
formés des droites menées et de droites données soit égale à une aire donnée,
le point est sur une droite donnée de position.\footnote{« Si sint quodcunq;
rectæ lineæ positione datæ, atq; ad ipsas a quodam puncto ducantur rectæ in
datis angulis, sit autem quod sub ductis et datis efficitur dato spatio æquale,
punctum rectam lineam positione datam continget » (vue 12). « Quod sub ductis
et datis efficitur » est lu ici comme la somme des rectangles, chacun fait
d'une droite menée et d'une droite donnée ; c'est une interprétation, que
demande la conclusion.} La démonstration moderne tient en une ligne. Une
droite menée du point $(x,y)$ à une droite donnée, sous un angle donné, est
proportionnelle à la distance du point à cette droite, donc, d'un même côté de
celle-ci, de la forme $d_i = \alpha_i x + \beta_i y + \gamma_i$. La condition
$\sum_i c_i d_i = K$ est du premier degré : c'est une droite. L'énoncé exact
demande deux précisions que le feuillet ne donne pas : les longueurs $d_i$
étant positives, le lieu est la partie de cette droite qui tombe dans la région
où chaque $d_i$ garde le signe choisi ; et si
$\sum c_i\alpha_i = \sum c_i\beta_i = 0$, la condition ne dépend plus du point,
et le lieu est toute la région ou rien.

\subsection*{Le rectangle des inconnues : l'hyperbole rapportée à ses asymptotes}

Le « second degré » des équations locales est $xy = Z$, avec $Z$ une aire. Le
point $I$ est sur une hyperbole : on mène par $N$ la parallèle $NR$ à $ZI$, on
choisit $M$ sur $NZM$ et $MO$ parallèle à $ZI$ de sorte que $NM\cdot MO = Z$, et
l'hyperbole qui passe par $O$ et a pour asymptotes $NR$ et $NM$ passe par $I$.
C'est la propriété de l'hyperbole rapportée à ses asymptotes : dans ce repère,
droit ou oblique, son équation est $xy = $ constante.

Toute équation où se mêlent le rectangle $xy$, des termes du premier degré et
une constante, sans carré, s'y ramène. L'exemple est $D + xy = Rx + Sy$, où $D$
est une aire. Comme
\[ (x - S)(R - y) = Rx + Sy - xy - RS, \]
l'équation s'écrit
\[ (x - S)(R - y) = D - RS , \]
ce que le feuillet dit en ôtant $RS$ de l'aire $D$. Le point $I$ est sur
l'hyperbole d'asymptotes les droites $x = S$ et $y = R$.\footnote{Le feuillet
suppose $D > RS$, puisqu'il construit un rectangle égal à $D - RS$. Si
$D < RS$, l'équation s'écrit $(x-S)(y-R) = RS - D$ : c'est encore une
hyperbole de mêmes asymptotes, dont on prend les deux autres quarts. Si
$D = RS$, le lieu dégénère en la paire de droites $x = S$ et $y = R$. En
général, $bxy + dx + ey + f = 0$ avec $b \neq 0$ s'écrit
$(bx + e)(by + d) = de - bf$ : une hyperbole d'asymptotes parallèles aux axes
si $de \neq bf$, deux droites sinon.} La construction : $NO = S$ sur $NZM$, $Nd$
parallèle à $ZI$ et égale à $R$, $dp$ parallèle à $NM$ et rencontrée en $p$ par
$ZI$ prolongée ; $u$ est sur $dp$ au-dessus de $O$. Alors $up = OZ = x - S$ et
$pI = Zp - ZI = R - y$, donc $up \cdot pI = D - RS$, et l'hyperbole
d'asymptotes $pu$ et $uO$ menée par un point $y$ tel que $ux\cdot xy = D - RS$
passe par $I$.\footnote{Le point que la vue 13 écrit « d » est celui que la vue
14 écrit « D » ; les lettres $x$ et $y$ de ce point de la figure sont celles
du feuillet, non les coordonnées. Le feuillet justifie $pI = R - y$ par « cum
uD siue ZI æquetur R » : tel qu'il est lu, $uD$ est un segment du côté
supérieur, égal à $NO = S$, et $ZI$ vaut $y$. Il faut un côté vertical du
rectangle, $Nd$, $Ou$ ou $Zp$, qui vaut $R$ ; la note de la marge, « + forsan
Zp », corrige le second terme, et cette lecture la suit.}

\subsection*{Les équations homogènes du second degré : des droites}

Le « degré suivant » est celui où $x^2$ est égal à $y^2$, ou dans un rapport
donné avec lui, ou encore $x^2 + xy$ dans un rapport donné avec $y^2$ : toutes
les équations dont chaque terme est de degré deux exactement. Le texte conclut
que $I$ est sur une droite : si $(x^2 + xy)/y^2$ est donné, et qu'on mène dans
le triangle $NZI$ une parallèle $OR$ à $ZI$, les triangles $NOR$ et $NZI$ sont
semblables, et $(NO^2 + NO\cdot OR)/OR^2$ a la même valeur ; la droite $NI$
convient.\footnote{« Iuncta NR dabitur positione \& satisfaciet proposito » ;
« satisfaciet » est d'une lecture un peu douteuse, la fin du mot serrée contre
la marge.} C'est juste, pour des grandeurs positives. L'énoncé moderne est un
peu plus riche : $ax^2 + bxy + cy^2 = 0$ se factorise en deux droites passant
par $N$ si $b^2 - 4ac > 0$, en une droite double si $b^2 - 4ac = 0$, et se
réduit au seul point $N$ si $b^2 - 4ac < 0$. Dans l'exemple, en posant
$t = x/y$, on a $t^2 + t - k = 0$, qui a une racine positive et une négative :
la première donne la demi-droite du feuillet, la seconde une autre droite par
$N$, qui passe hors de l'angle où l'écrivain place ses grandeurs. Ces paires de
droites sont les coniques dégénérées de la classification.

\section*{L'Isagoge, suite : parabole, cercle, ellipse, hyperbole (ff. 4v–7v, vues 16 à 22)}

\pagerange{16}{22}
« Si aux puissances des inconnues ou à leurs rectangles se mêlent des termes
en partie entièrement donnés, en partie faits d'une droite donnée et de l'une
des inconnues, la construction sera plus difficile. Nous construisons
brièvement chaque cas, et nous le démontrons. » Ce sont les coniques propres.

\subsection*{La parabole}

Si $x^2 = Dy$, le point $I$ est sur une parabole de sommet $N$, de diamètre la
parallèle $NP$ à $ZI$ et de paramètre $D$ : par construction,
$D\cdot NP = PI^2$, c'est-à-dire $Dy = x^2$.\footnote{Le feuillet écrit
« applicatæ sint parallelæ NP » ; les ordonnées d'une parabole de diamètre $NP$
sont parallèles à $NZ$, non à $NP$, et c'est ce que demande la suite. Il écrit
ensuite « rectangulum sub D in IP æquabitur quadrato NZ », où il faut $NP$ (ou
$ZI$), puisque $IP = NZ = x$ ; la conclusion qu'il en tire, « D in E æquabitur
Aq », est la bonne. Il écrit « Parobolen », et biffe un « erit » en fin de
ligne pour le récrire à la suivante. Dans « quadrato N$\overset{Z}{E}$ », un Z
est écrit au-dessus du E, de la même encre : la leçon corrigée est « quadrato
NZ ».} Si l'angle $\widehat{NZI}$ est droit, $NP$ est l'axe ; sinon c'est un
diamètre, et c'est bien le mot du texte.

« À cette équation se ramèneront très facilement toutes celles où $x^2$ se mêle
à des termes faits de grandeurs données et de $y$, ou $y^2$ à des termes faits
de grandeurs données et de $x$, et de même s'il s'y mêle des termes entièrement
donnés. » C'est la complétion du carré : $x^2 + ax + by + c = 0$ s'écrit
$\bigl(x + \tfrac a2\bigr)^2 = -b\,\bigl(y + \tfrac{c}{b} - \tfrac{a^2}{4b}\bigr)$,
une parabole, pourvu que $b \neq 0$ ; si $b = 0$, il reste deux droites
parallèles, une droite ou rien. La condition $b\neq 0$ est celle que l'énoncé
du feuillet contient en exigeant des termes « sub datis in E ».

Le feuillet traite ensuite trois cas particuliers :
\begin{itemize}
  \item $y^2 = Dx$ : parabole de sommet $N$, de diamètre $NZ$, de paramètre $D$,
    les ordonnées parallèles à $NP$ ;
  \item $B^2 - x^2 = Dy$ : si $DR = B^2$, l'équation devient
    $D(R - y) = x^2$ ; on porte $NM = R$ parallèlement à $ZI$, on mène $MO$
    parallèle à $NZ$, et $OI = R - y$, $MO = x$ : parabole de sommet $M$, de
    diamètre $MN$, de paramètre $D$, tournée vers $N$ ;
  \item $B^2 + x^2 = Dy$ : $D(y - R) = x^2$, « et le reste comme plus haut ».
\end{itemize}
Tous trois sont justes.

\subsection*{Le cercle}

Si $B^2 - x^2 = y^2$ et que l'angle en $Z$ est droit, $I$ est sur le cercle de
centre $N$ et de rayon $B$. L'écrivain précise que s'y ramènent toutes les
équations en $x^2$, $y^2$ et termes du premier degré, « pourvu que l'angle
$\widehat{NZI}$ soit droit et que les coefficients de $x^2$ égalent ceux de
$y^2$ ». Son exemple :
\[ B^2 - 2Dx - x^2 = y^2 + 2Ry . \]
On ajoute $R^2$, puis $D^2$, aux deux membres, et l'on pose
$P^2 = B^2 + R^2 + D^2$ :
\[ P^2 - (x + D)^2 = (y + R)^2 , \]
de sorte que, en prenant $x + D$ et $y + R$ pour nouvelles inconnues, on
retombe sur le cas précédent : un cercle de centre le point $(-D, -R)$ et de
rayon $P$.\footnote{Le « $+$ D » de « loco ipsius A $+$ D » est écrit
au-dessus de la ligne ; un second « $+$ D », très pâle, se voit au-dessus du
second A. Le feuillet écrit aussi « Pq $-$ Dq $-$ D in A$''$ $-$ Aq æquabitur
Rq $+$ Bq $-$ D in A$''$ $-$ Aq », qui n'est que la reformulation de
$P^2 - D^2 = R^2 + B^2$. En général, $x^2 + y^2 + ax + by + c = 0$ est un
cercle si $\tfrac{a^2}{4} + \tfrac{b^2}{4} - c > 0$, un point s'il est nul, rien
s'il est négatif ; dans l'exemple, $P^2 = B^2 + D^2 + R^2$ est toujours
positif.}

L'écrivain ajoute : « c'est par cette voie que nous avons construit toutes les
propositions du second livre des Lieux plans d'Apollonius, et démontré que les
six premières ont lieu pour des points quelconques, ce qui est assurément
admirable, et peut-être ignoré d'Apollonius ». Cette lecture n'en peut rien
contrôler : les énoncés d'Apollonius ne sont pas sur les feuillets.\footnote{Le
dernier mot de la vue 18, biffé d'un trait épais, se lit peut-être
« quando » ; il n'a pas de suite.}

\subsection*{L'ellipse}

Si $(B^2 - x^2)/y^2$ est un rapport donné, $I$ est sur une ellipse de centre
$N$, qui passe par le point $M$ tel que $NM = B$, de diamètre $NM$, les
ordonnées parallèles à $ZI$. Car $B^2 - x^2 = (B - x)(B + x)$ est le produit
des deux segments du diamètre, et une ellipse est la courbe où le carré de
l'ordonnée est dans un rapport constant avec ce produit. Aujourd'hui :
$y^2 = k(B^2 - x^2)$, soit $\dfrac{x^2}{B^2} + \dfrac{y^2}{kB^2} = 1$.

S'y ramènent, dit le texte, les équations où $x^2$ et $y^2$ sont de part et
d'autre de l'égalité avec des signes contraires et des coefficients
différents — autrement dit, écrites d'un même côté, avec le même signe. Si les
coefficients sont égaux et l'angle droit, c'est un cercle ; « si les
coefficients sont les mêmes, mais que l'angle ne soit pas droit, le lieu sera à
une ellipse ». Tout cela est exact : $x^2 + y^2 = B^2$ en coordonnées obliques
est une ellipse, image d'un cercle par une application affine, et non un
cercle ; et le cercle est la seule des coniques qui dépende de l'angle.
L'équation $ax^2 + cy^2 + dx + ey + f = 0$ avec $ac > 0$ est une ellipse, ou un
point, ou rien, selon le signe de ce qui reste après complétion des carrés ;
les deux derniers cas, le feuillet ne les envisage pas.\footnote{« Et licet
immisceantur æquationibus homogenea sub datis \& A vel E fiet reductio eo quod
iam \emph{vsurpauimus} artificio » (vue 21) : la réduction des termes du
premier degré se fait par l'artifice déjà employé, la complétion des carrés.
« vsurpauimus » est d'une lecture douteuse.}

\subsection*{L'hyperbole}

Si $(x^2 + B^2)/y^2$ est un rapport donné, $I$ est sur une hyperbole.\footnote{Le
mot « Hyperbolen » de la vue 21 est lu avec doute, la fin étant une boucle qui
peut être une abréviation ; c'est le mot que demandent la construction et la
conclusion. Un petit « m » surélevé, après la parenthèse « (posita MN æquali
NR) », n'est pas compris.} Écrivons le rapport donné sous la forme $B^2/NR^2$,
ce qui fixe un point $R$ sur la parallèle $NO$ à $ZI$ menée par $N$ ; soit $M$
le symétrique de $R$ par rapport à $N$. L'hyperbole de centre $N$, de sommet
$R$, de diamètre $RO$, les ordonnées parallèles à $NZ$, est celle où
\[ \frac{MO\cdot OR}{OI^2} = \frac{NR^2}{B^2} . \]
Le texte ajoute $NR^2$ et $B^2$ aux deux termes du premier rapport
(« componendo ») : comme $MO\cdot OR + NR^2 = NO^2$ et que $NO = ZI = y$,
$OI = NZ = x$, il vient
\[ \frac{y^2}{x^2 + B^2} = \frac{NR^2}{B^2} ,
   \qquad\text{soit}\qquad \frac{y^2}{NR^2} - \frac{x^2}{B^2} = 1 , \]
qui est l'équation cherchée.\footnote{Le texte conclut « Punctum igitur Z est
ad Hyperbolen » ; c'est le point $I$, celui de l'énoncé de la vue 21. La note
marginale de la vue 22, appelée par un astérisque et de la même main, complète
le raisonnement : « æquatur quadrato NZ siue Aq vnâ cum Bq ».} Et le texte :
s'y ramènent toutes les équations en $x^2$ et $y^2$ « pourvu que $x^2$ ait, de
l'autre côté, le même signe que $y^2$ ; car si les signes sont différents, la
proposition se conclut par des cercles ou des ellipses ». C'est la règle
exacte : écrites d'un même côté, $ax^2 + cy^2 + \dots$ avec $ac < 0$ donne une
hyperbole — ou, si la constante qui reste après complétion des carrés est
nulle, la paire de ses asymptotes, cas que le feuillet n'examine pas.

\section*{L'Isagoge, fin : le terme rectangle, le bilan, la proposition finale (ff. 7v–9v, vues 22 à 26)}

\pagerange{22}{26}
\subsection*{« La plus difficile de toutes les équations »}

Reste le cas où le rectangle $xy$ se mêle aux carrés. L'exemple :
\[ B^2 - 2x^2 = 2xy + y^2 . \]
En ajoutant $x^2$ aux deux membres,
\[ B^2 - x^2 = x^2 + 2xy + y^2 = (x + y)^2 . \]
Si l'on prenait $x + y$ pour seconde coordonnée, ce serait le cercle de centre
$N$ et de rayon $B$ ; mais c'est l'extrémité $V$ de $y$ qu'on cherche. Le texte
porte donc $ZI = x + y$ (le point $I$ décrit ce cercle) et, sur $ZI$, le point
$V$ tel que $VI = NZ = x$, de sorte que $ZV = y$. Il mène $MR$ parallèle à $ZI$
et égale à $MN$, joint $NR$, que $IZ$ prolongée coupe en $O$. Le triangle $NMR$
étant isocèle, $ZO = NZ = x$, donc $VO = VZ + ZO = x + y = ZI$, et
\[ OV^2 = MN^2 - NZ^2 . \]
Le triangle $NMR$ étant de forme donnée, $NZ^2/NO^2 = NM^2/NR^2$ est donné, donc
aussi $(MN^2 - NZ^2)/(NR^2 - NO^2)$, et par suite $(NR^2 - NO^2)/OV^2$. Les
points $N$ et $R$ sont donnés, l'angle $\widehat{NOZ}$ aussi : par le cas de
l'ellipse, $V$ est sur une ellipse de centre $N$, de diamètre porté par $NR$,
les ordonnées parallèles à $ZI$.\footnote{À la vue 24, une encre plus noire
récrit « NZ » au-dessus de « VZ » et « NR » au-dessus de « VR », sans barrer les
leçons de la ligne. Le raisonnement demande $NZ$ et $NR$ ; cette lecture suit
la correction.} La vérification moderne : avec l'angle droit, $NO = x\sqrt2$,
$NR = B\sqrt2$ et $OV = x + y$ ; l'équation $2x^2 + 2xy + y^2 = B^2$ s'écrit
$NR^2 - NO^2 = 2\,OV^2$, et, dans le repère oblique d'axes $NR$ et $ZI$, le
point $V$ est sur l'ellipse $\dfrac{X^2}{2B^2} + \dfrac{Y^2}{B^2} = 1$. Le
discriminant de $2x^2 + 2xy + y^2$ est $2^2 - 4\cdot2\cdot1 = -4 < 0$ : c'est
bien une ellipse.

« Par une méthode semblable se ramènent aux cas précédents tous les autres
[…] ; car c'est toujours au moyen d'un triangle de forme connue que la question
se construit. » Le triangle de forme donnée est un changement de coordonnées
obliques qui fait disparaître le rectangle. En termes modernes, si $c \neq 0$,
\[ ax^2 + bxy + cy^2
   = c\Bigl(y + \frac{b}{2c}\,x\Bigr)^2 + \frac{4ac - b^2}{4c}\,x^2 ; \]
prendre $y + \tfrac{b}{2c}x$ pour nouvelle coordonnée revient à mesurer
l'ordonnée le long d'une autre direction, et la conique se trouve ramenée à un
cas sans rectangle. Les deux nouveaux coefficients des carrés, $c$ et
$(4ac - b^2)/(4c)$, ont pour produit $(4ac - b^2)/4$ : de même signe, c'est une
ellipse ; de signes contraires, une hyperbole ; si l'un est nul, une parabole
ou une conique dégénérée. Le type se lit sur le signe de $b^2 - 4ac$ : c'est le
discriminant.\footnote{Le cas $c = 0$, $a \neq 0$ se
traite en échangeant les rôles de $x$ et de $y$ ; le cas $a = c = 0$ est celui
de l'hyperbole rapportée à ses asymptotes, vu plus haut.} L'écrivain conclut :
« Breuiter igitur \& dilucidè complexi sumus quidquid de Locis planis \&
Solidis inexplicatum veteres reliquere. »\footnote{« Nous avons donc embrassé,
brièvement et clairement, tout ce que les Anciens ont laissé inexpliqué sur les
lieux plans et solides. » Il ajoute qu'on saura désormais à quel lieu
appartient chaque cas de la dernière proposition du livre I des Lieux plans
d'Apollonius, que cette lecture ne peut pas contrôler.} La prétention est
juste pour les coniques propres : toute équation du second degré définit une
conique ou une conique dégénérée, et la réduction ci-dessus le démontre. Le
feuillet, lui, ne démontre pas que ses réductions épuisent tous les cas, et ne
nomme pas les cas dégénérés.

\subsection*{Le bilan, dans la classification d'aujourd'hui}

Pour l'équation $ax^2 + bxy + cy^2 + dx + ey + f = 0$, voici les cas de
l'Isagoge et la classification par le discriminant $\Delta = b^2 - 4ac$, là où
les deux coïncident :
\begin{itemize}
  \item $a = b = c = 0$ : une droite (vues 10 à 12) ;
  \item $a = c = 0$, $b \neq 0$, donc $\Delta > 0$ : une hyperbole
    d'asymptotes parallèles aux axes, ou deux droites (vues 12 à 14) ;
  \item $d = e = f = 0$ : une ou deux droites par l'origine, ou l'origine
    seule (vues 14 et 15) ;
  \item un seul carré, $b = 0$, donc $\Delta = 0$ : une parabole si le terme du
    premier degré dans l'autre inconnue est présent (vues 16 à 18) ;
  \item $b = 0$, $ac > 0$, donc $\Delta < 0$ : une ellipse, un cercle si $a = c$
    et l'angle droit, ou un point, ou rien (vues 18 à 21) ;
  \item $b = 0$, $ac < 0$, donc $\Delta > 0$ : une hyperbole, ou deux droites
    sécantes (vues 21 et 22) ;
  \item $b \neq 0$ et un carré au moins : réduit par un triangle de forme
    donnée à l'un des cas précédents ; l'exemple a $\Delta < 0$ et donne une
    ellipse (vues 22 à 24).
\end{itemize}

\subsection*{La proposition finale}

« En guise de couronnement », l'écrivain ajoute une proposition « très belle,
dont la facilité apparaîtra aussitôt » : si, d'un même point, on mène à des
droites données de position, en nombre quelconque, des droites sous des angles
donnés, et que les carrés de toutes les droites menées soient ensemble égaux à
une aire donnée, le point est sur un lieu solide donné de position.\footnote{« Et
sint species ab omnibus ductis dato spatio æqualis » : « species » est lu ici
comme les carrés des droites menées, pris ensemble ; c'est une interprétation,
que demande le mot de « lieu solide ».} C'est vrai, et l'on peut dire mieux.
Chaque droite menée est une fonction affine $d_i = \alpha_i x + \beta_i y +
\gamma_i$ du point, et la condition $\sum_i d_i^2 = K$ est du second degré ; sa
partie quadratique $\sum_i(\alpha_i x + \beta_i y)^2$ est positive, et même
définie positive dès que deux des droites données ne sont pas parallèles. Le
lieu est alors une ellipse (un cercle dans certains cas), réduite à un point ou
vide si $K$ est trop petit ; si toutes les droites sont parallèles, ce sont deux
droites parallèles, une droite ou rien.

L'exemple qui suit n'est pas un cas de cet énoncé : il porte sur des distances
à deux points, et sur un rapport à un triangle, non sur des droites. Il
illustre la méthode. Étant donnés deux points $N$ et $M$, trouver le lieu des
points $I$ tels que $(IN^2 + IM^2)$ soit dans un rapport donné avec le triangle
$INM$. Avec $NM = B$, $Z$ le pied de la perpendiculaire abaissée de $I$,
$NZ = x$ et $ZI = y$ :
\[ IN^2 + IM^2 = x^2 + y^2 + (B - x)^2 + y^2 = 2x^2 + B^2 - 2Bx + 2y^2 , \]
que le texte met en rapport avec le rectangle $By$, double du triangle ; le
rapport reste donné. La construction du feuillet (vue 26) : couper $NM$ en deux
en $Z$, élever la perpendiculaire $ZV$ telle que le rapport donné soit celui de
$4\,ZV$ à $ZM$ ; sur $VZ$ comme diamètre, décrire un demi-cercle et y inscrire
la corde $ZO = ZM$ ; le cercle de centre $V$ et de rayon $VO$ est le lieu : pour
tout point $R$ de ce cercle, $(RN^2 + RM^2)/\text{triangle }RNM$ est le rapport
donné.\footnote{La ligne de la construction est corrigée à plusieurs endroits
(« ipsi » et « ZO » barrés, « ZM » écrit puis barré) ; la leçon finale, que
suit cette lecture, est « applicetur ZO æqualis ipsi ZM ». Un grand cercle très
pâle, au crayon, autour de la figure de la vue 25, n'est pas de l'encre du
texte.}

La vérification. Notons $\rho = 4\,ZV/ZM$ le rapport donné et $h = ZV$. Dans le
demi-cercle, l'angle $\widehat{VOZ}$ est droit, donc $VO^2 = h^2 - ZM^2$. Si
$R = (x,y)$ est sur le cercle de centre $V = (B/2, h)$ et de rayon $VO$,
\[ \Bigl(x - \frac B2\Bigr)^2 + y^2 = 2hy - ZM^2 , \]
d'où, puisque $ZM = B/2$,
\[ RN^2 + RM^2 = 2\Bigl(x - \frac B2\Bigr)^2 + 2y^2 + \frac{B^2}{2} = 4hy , \]
tandis que le triangle vaut $ZM\cdot y$ : le rapport est $4h/ZM = \rho$. La
construction est juste, sous une condition qu'elle cache : la corde $ZO = ZM$
ne s'inscrit dans le demi-cercle que si $ZM < ZV$, c'est-à-dire $\rho > 4$. Et
en effet $IN^2 + IM^2 \geq 2\,IN\cdot IM \geq 4\cdot\text{triangle }INM$, avec
égalité seulement quand le triangle est rectangle isocèle en $I$ : pour
$\rho = 4$ le lieu est le seul point $V$, pour $\rho < 4$ il est vide. Enfin,
l'aire du triangle ne dépendant pas du côté de $NM$ où se trouve $I$, le lieu
complet est formé de ce cercle et de son symétrique par rapport à $NM$. L'Isagoge
s'arrête là, au milieu de la vue 26.

\section*{« Methodus ad disquirendam maximam \& minimam » (f. 10r–v, vues 27 et 28)}

\pagerange{27}{28}
Un nouveau feuillet commence par un titre français : « Copie d'un escrit enuoyé
par le R. Pere Marçenne a des Cartes ».\footnote{« Monsieur », entre « a » et
« des Cartes », est biffé et d'une lecture douteuse. Ce titre nomme l'envoyeur
et le destinataire d'un écrit, le P. Mersenne et Descartes ; il ne nomme pas
son auteur, et la pièce non plus.} Suit le texte latin, qui tient en une
règle et un exemple.

\subsection*{La règle}

« Toute la doctrine de l'invention du maximum et du minimum repose sur deux
positions inconnues et sur cette seule règle. » Qu'on prenne pour inconnue
$A$ ce qu'il convient dans la question, et qu'on écrive la grandeur à rendre
maximale ou minimale en fonction de $A$ ; puis qu'on remplace $A$ par $A + E$
et qu'on écrive de nouveau la grandeur. « Adæquentur, vt Loquitur Diophantus » :
qu'on « adégale », comme dit Diophante, les deux expressions ; qu'on ôte les
termes communs ; qu'on divise tout par $E$, ou par une puissance de $E$, jusqu'à
ce que l'un des termes soit libre de $E$ ; qu'on efface alors les termes qui
contiennent encore $E$, et qu'on égale le reste.\footnote{« Adæqualitas »,
« adéquation », rend chez les traducteurs latins de Diophante un terme grec qui
signifie « presque-égalité » ; c'est un renseignement général, que les feuillets
ne donnent pas. Le feuillet emploie « adæquentur », « adæquabitur » pour cette
presque-égalité et « æquabitur » pour l'égalité finale ; cette lecture écrit
la première $\approx$ et la seconde $=$. Si, après la suppression des termes
communs, l'un des membres est vide, le feuillet égale « les termes niés aux
termes affirmés », ce qui revient au même.} Ainsi, en écrivant $x$ pour $A$,
$e$ pour $E$ et $f$ pour la grandeur :
\[ f(x + e) \approx f(x), \qquad
   \frac{f(x+e) - f(x)}{e} \approx 0, \qquad
   \text{puis } e \text{ effacé}. \]

L'exemple : partager une droite $AC = B$ en un point $E$ de sorte que le
rectangle $AE\cdot EC$ soit maximal. Avec $AE = x$, le rectangle vaut
$f(x) = Bx - x^2$, et
\[ f(x + e) = Bx - x^2 + Be - 2xe - e^2 . \]
Adégalant et ôtant les termes communs, $Be \approx 2xe + e^2$ ; divisant par
$e$, $B \approx 2x + e$ ; effaçant $e$, $B = 2x$. Il faut couper la droite en
son milieu. Le feuillet conclut : « Nec potest generalior dari
methodus. »\footnote{« On ne peut donner de méthode plus générale. »}

\subsection*{Ce que calcule la règle, et ce qu'il faut pour qu'elle démontre}

Pour une grandeur $f$ polynomiale en $x$, la différence $f(x+e) - f(x)$ est un
polynôme en $e$ sans terme constant ; elle s'écrit exactement
\[ f(x+e) - f(x) = e\,\bigl(f'(x) + e\,g(x,e)\bigr), \]
où $f'$ est la dérivée et $g$ un polynôme. Diviser par $e$ est donc légitime
($e \neq 0$), et « effacer $e$ » dans le quotient, c'est prendre sa valeur en
$e = 0$, qui est sa limite quand $e \to 0$ : la règle calcule $f'(x)$ et écrit
$f'(x) = 0$. Elle donne les points stationnaires de $f$.

Ce n'est pas, telle qu'elle est écrite, une démonstration. L'« adéquation »
n'est pas une égalité — en un maximum, $f(x+e)$ diffère en général de $f(x)$
pour $e \neq 0$ —, et l'on ne peut à la fois diviser par $e$, qui suppose
$e \neq 0$, et poser ensuite $e = 0$ sans dire ce qui l'autorise. Ce qui en
fait une démonstration aujourd'hui, c'est le théorème suivant : si $f$ est
dérivable en un point $x$ intérieur à son domaine et y atteint un maximum ou un
minimum local, alors $f'(x) = 0$. Pour un polynôme la preuve est dans la
formule ci-dessus : si $f'(x) \neq 0$, le facteur $f'(x) + e\,g(x,e)$ garde le
signe de $f'(x)$ pour $e$ assez petit, et $f(x+e) - f(x)$ change de signe avec
$e$ ; $x$ n'est donc pas un extremum.\footnote{On appelle aujourd'hui ce
résultat « théorème de Fermat » sur les extremums intérieurs ; le nom est une
attribution moderne. Les feuillets n'emploient ni la dérivée, ni la limite, ni
aucun énoncé de ce genre ; la dérivée, sous ce nom et dans ce formalisme, est
postérieure à Mersenne et à Descartes, que nomme le titre.} La règle donne donc
une condition nécessaire, et seulement pour un extremum intérieur : elle ne
distingue pas le maximum du minimum, elle retient aussi des points qui ne sont
ni l'un ni l'autre ($x^3$ en $0$), et elle ne voit pas les extrémités d'un
intervalle.

Dans l'exemple, la conclusion est juste et se démontre en une ligne :
$f\bigl(\tfrac B2 + h\bigr) = \tfrac{B^2}{4} - h^2 \leq f\bigl(\tfrac B2\bigr)$,
avec égalité seulement pour $h = 0$ ; le milieu donne le maximum strict.

\section*{« De tangentibus linearum curvarum » (f. 11r–v, vues 29 et 30)}

\pagerange{29}{30}
« Nous ramenons à la méthode précédente l'invention des tangentes en des
points donnés de lignes courbes quelconques. » La pièce est sur le feuillet qui
suit, de la même main ; qu'elle appartienne à l'« escrit » que désigne le titre
de la vue 27, le feuillet ne le dit pas, mais elle en continue la méthode.

\subsection*{La tangente à la parabole}

Soit une parabole $Bdn$ de sommet $d$ et de diamètre $dc$, un point $B$ de la
courbe, et la tangente $BE$ cherchée, qui coupe le diamètre en $E$. On prend
sur $BE$ un point $O$, et l'on mène les ordonnées $OI$ et $BC$. Comme $O$ est
hors de la parabole,
\[ \frac{cd}{di} > \frac{BC^2}{OI^2} , \]
et, par les triangles semblables, $BC^2/OI^2 = CE^2/IE^2$ ; donc
$cd/di > CE^2/IE^2$. Posons $cd = d$, donnée,\footnote{Le texte écrit d'abord
« Sit igitur cd æqualis B data », puis se sert de la lettre d ; la note
marginale française, appelée par un astérisque et qui parle de l'auteur à la
troisième personne, le relève : « Il a icy nommé B ce qu'il nomme d par
apres ».} $CE = a$ l'inconnue (le feuillet l'appelle $A$) et $CI = e$ (le
feuillet : $E$). Alors $di = d - e$ et $IE = a - e$, et
\[ \frac{d}{d - e} > \frac{a^2}{a^2 + e^2 - 2ae} ,
   \qquad\text{soit}\qquad
   d a^2 + d e^2 - 2dae > d a^2 - a^2 e . \]
On « adégale » : ôtant $da^2$, $de^2 + a^2 e \approx 2dae$ ; divisant par $e$,
$de + a^2 \approx 2da$ ; effaçant $de$, $a^2 = 2da$, d'où $a = 2d$. « Nous avons
prouvé que $CE$ est double de $cd$, ce qui est bien le cas. » C'est exact : pour
la parabole $y^2 = px$ rapportée à son sommet et à son axe, la tangente au
point d'abscisse $x_0$ coupe l'axe à la distance $2x_0$ de ce point ; la
sous-tangente est double de l'abscisse.\footnote{Aux vues 28 et 30, un petit rond
ou un petit trait posé sur un signe moins ne le change pas : les lignes
suivantes le confirment.}

\subsection*{Ici, la règle est une démonstration}

L'inégalité du feuillet est exacte, et c'est elle qui fait de ce calcul une
preuve. Posons
\[ F(e) = d\,(a - e)^2 - (d - e)\,a^2 = e\,\bigl(a^2 - 2da + de\bigr) . \]
Dire que la droite $BE$ touche la parabole en $B$ sans la couper, c'est dire
que ses points autres que $B$ sont hors de la courbe, c'est-à-dire $F(e) > 0$
pour tout $e \neq 0$ voisin de $0$, des deux côtés de $B$ ; et $F(0) = 0$. Le
point $e = 0$ est donc un minimum de $F$, et le théorème de la section
précédente donne $F'(0) = a^2 - 2da = 0$ : c'est exactement l'« adéquation ».
Réciproquement, pour $a = 2d$, $F(e) = d\,e^2 > 0$ pour $e \neq 0$ : la droite
reste hors de la parabole, c'est la tangente. Le feuillet ne considère que des
points $O$ entre $B$ et $E$, c'est-à-dire $e > 0$ ; c'est l'inégalité des deux
côtés qu'il faut, et elle est vraie.

\subsection*{Les annonces}

« Et la méthode ne trompe jamais. Bien plus, elle peut s'étendre à la plupart
des plus belles questions » : les centres de gravité des figures limitées par
des courbes et des droites, et des solides ; les quadratures, et le rapport des
solides qu'elles engendrent aux cônes de même base et de même hauteur, « dont
nous avons déjà traité abondamment avec M. de Roberval ». Ces annonces sont
rapportées ; le feuillet n'en donne aucun calcul.\footnote{« Nec vnquam fallit
methodus » est trop dire, tel quel : pour les maxima, la règle retient aussi
des points stationnaires qui ne sont pas des extremums ; pour les tangentes,
elle suppose la courbe régulière au point considéré, et, en un point double,
elle ne donne rien. Les exemples du feuillet ne rencontrent aucun de ces cas.
Le texte s'arrête sur « egimus », sans point.}

\section*{« Appendix ad Isagogen topicam » : les problèmes solides résolus par les lieux (ff. 12r–14r, vues 31 à 35)}

\pagerange{31}{35}
Le titre complet est « Appendix ad Isagogen topicam continens solutionem
problematum solidorum per locos » : « appendice à l'introduction aux lieux,
contenant la solution des problèmes solides par les lieux ». Le texte commence
par « Patuit methodus quâ lineæ locales deteguntur » et renvoie plusieurs fois à
« notre méthode » : il se donne pour la suite d'une introduction aux lieux.
Que cette introduction soit l'Isagoge des vues 9 à 26 est le rapprochement
naturel, et c'est celui de cette lecture ; le feuillet ne le dit pas en toutes
lettres, la main est autre, et les nombres 43 à 45 du pied des pages
appartiennent à un autre assemblage.

\subsection*{Le principe}

Un problème « solide » conduit à une équation du troisième ou du quatrième
degré en une inconnue. « Il faut restreindre la liberté qu'ont les inconnues
d'errer hors de leurs limites » : une équation à deux inconnues a une infinité
de solutions, qui forment une ligne ; deux telles équations, deux lignes, se
coupent, et le point d'intersection ramène la question « de l'infini à des
termes prescrits ».\footnote{La phrase est corrigée en plusieurs endroits sur le
feuillet ; la leçon voulue paraît « Secant quippe se inuicem duæ lineæ locales
positione datæ, et punctum intersectionis positione datum quæstionem ex
Infinito ad terminos præscriptos adigit », avec un mot barré illisible après
« determinatur » et un autre devant « Sectionis ».} Aujourd'hui : on introduit
une seconde inconnue $y$ de sorte que l'équation se coupe en deux équations en
$(x,y)$, chacune une conique ; les abscisses des points communs sont les
racines.

\subsection*{La cubique par une parabole et une hyperbole}

Soit $x^3 + Bx^2 = ZB$, où $Z$ est une aire. On égale chacun des deux membres
au solide $Bxy$. D'un côté, $x^3 + Bx^2 = Bxy$, soit, en divisant par $x$,
$x^2 + Bx = By$ : une parabole. De l'autre, $ZB = Bxy$, soit $xy = Z$ : une
hyperbole. Le point commun donne $x$. En effet, sur l'hyperbole $y = Z/x$, et
l'équation de la parabole redonne $x^3 + Bx^2 = BZ$. Pour $x > 0$ il y a un
seul point commun — la parabole croît à partir de $0$, l'hyperbole décroît
depuis l'infini —, et c'est l'unique racine positive ; les autres points
communs, s'il y en a, ont $x < 0$ et donnent les racines négatives, dont le
feuillet, qui ne connaît que des grandeurs positives, ne parle
pas.\footnote{Le premier « æquari » de l'énoncé est biffé en fin de ligne et
récrit à la suivante.} Il en va de même, dit le texte, de toute cubique : pour
$x^3 + ax^2 + bx = c$, on égale les deux membres à $kxy$, et il vient
$x^2 + ax + b = ky$ (une parabole) et $c = kxy$ (une hyperbole).

\subsection*{La quartique par une parabole et un cercle}

Soit $x^4 + Bx + Z^2x^2 = D$ ($B$ est un volume, $D$ une grandeur de dimension
quatre).\footnote{Le feuillet écrit le coefficient de $x^2$ « Zpl. », un plan,
puis « Zq », un carré : il prend ce plan pour le carré d'une droite $Z$, ce
qui est toujours possible. Les Z de ces pages sont d'une encre plus noire,
« comme remplis après coup ».} Donc $x^4 = D - Bx - Z^2x^2$, et l'on égale les
deux membres à $Z^2y^2$.\footnote{Le feuillet n'écrit aucun signe entre
« Dppl. », « B Sol. in A » et « Zq in Aq » (vue 31), ni entre « Dppl. » et
« B Solid. » (vue 32) ; les moins sont ceux que demande l'équation de départ, et
la suite du feuillet les a (« $-$ Zq in Aq », « $-$ B Sol. in A/Zq »). Un
« Eq » isolé est barré sous la fraction.} D'un côté, $x^4 = Z^2y^2$ donne, par
« subdivision quadratique » — en prenant la racine carrée —, $x^2 = Zy$ : une
parabole. De l'autre, en divisant par $Z^2$,
\[ \frac{D}{Z^2} - \frac{B}{Z^2}\,x - x^2 = y^2 : \]
un cercle, de centre $\bigl(-\tfrac{B}{2Z^2}, 0\bigr)$ et de rayon
$\sqrt{D/Z^2 + B^2/(4Z^4)}$. Le calcul est juste : un point commun vérifie
$x^4 = Z^2y^2 = D - Bx - Z^2x^2$, et toute racine réelle $x$ donne le point
commun $(x, x^2/Z)$.

« Non dissimili methodo soluentur quæstiones omnes quadrato quadraticæ » : on
fait d'abord disparaître le terme en $x^3$ par la méthode de Viète, que le
feuillet cite (« cap. i. de emend. ») — aujourd'hui, la translation qui rend
la quartique réduite —, puis on met $x^4$ d'un côté, le reste de l'autre, et
l'on résout « par la parabole, le cercle ou l'hyperbole ».

\subsection*{Les deux moyennes proportionnelles}

Entre deux droites $B$ (la plus grande) et $D$, trouver deux moyennes
proportionnelles $x$ et $y$, c'est-à-dire $B/x = x/y = y/D$. La plus grande
vérifie $x^3 = B^2D$.\footnote{« Sint duæ rectæ B maior D minor » : « maior »
est lu d'après la suite, le mot étant tracé presque comme « minor » ; les
derniers mots de la ligne suivante, « nempe quod », sont d'une lecture
douteuse. L'astérisque de « major mediarum ponatur A » est dans la marge.} On
égale les deux membres à $Bxy$ : d'un côté $x^2 = By$, une parabole ; de l'autre
$xy = BD$, une hyperbole. La construction : sur une droite $OVN$, $OV = x$ et
la perpendiculaire $VM = y$ ; la parabole de sommet $O$, de paramètre $B$, de
diamètre parallèle à $VM$ ; l'hyperbole d'asymptotes $RO$ et $OV$ passant par le
point $Z$ tel que $On\cdot nZ = BD$ ; elles se coupent en $m$, et la
perpendiculaire $mV$ donne $OV$, la plus grande des deux
moyennes.\footnote{Le feuillet écrit « perpendicularis VZ » là où sa figure élève
la perpendiculaire en $n$ ; la lettre $V$ est le petit u du copiste, que la
transcription distingue mal du n. Il faut $nZ$. Les fins de « hyperboles » et
« parabolæ » sont sous un pâté ou surchargées.} La construction est juste, et
elle donne davantage que ce qu'elle dit : puisque $x^2 = By$, on a $B/x = x/y$,
et $VM = y$ est l'autre moyenne ; puisque $xy = BD$, $x/D = B/y$, ce qui
complète la proportion.

En multipliant par $x$, $x^4 = B^2Dx$ ; on égale les deux membres à $B^2y^2$ :
$x^2 = By$ et $Dx = y^2$, deux paraboles, dont l'intersection donne encore
$x$. C'est juste, à ceci près que les deux paraboles se coupent aussi à
l'origine, point introduit par la multiplication par $x$ et qu'il faut écarter ;
le feuillet ne le dit pas. « Prior constructio et posterior sunt apud Eutocium
in Archimedem » : les deux constructions sont chez Eutoce, dans son commentaire
d'Archimède, et se prêtent très facilement à cette méthode.\footnote{Le feuillet
ne nomme qu'Eutoce. C'est un renseignement général, non vérifié ici sur aucune
édition, que ce commentaire rapporte deux solutions attribuées à Ménechme, l'une
par une parabole et une hyperbole, l'autre par deux paraboles, qui sont
celles-ci.}

Suit une charge contre Viète : « Abeant igitur illæ parapleroses Vietææ », que
s'en aillent donc les « paraplérôses » de Viète, par lesquelles il ramène les
quartiques aux quadratiques par l'intermédiaire de cubiques « à racine plane » ;
quartiques et cubiques se résolvent ici avec la même élégance, « et l'on ne
peut, je crois, plus élégamment ». Ce que le texte décrit est la méthode de la
cubique résolvante : une quartique ramenée, par une cubique auxiliaire dont
l'inconnue est une aire, à deux équations du second degré.

\subsection*{Toutes les cubiques et quartiques par une parabole et un cercle}

« Pour que paraisse l'élégance de cette méthode » : la construction de tous les
problèmes cubiques et quartiques par la parabole et le cercle. Soit
$x^4 + Zx = D$ ($Z$ un volume, $D$ de dimension quatre), donc
$x^4 = D - Zx$.\footnote{Le feuillet écrit ici « ergo Aqq æquabitur Z Sol. in A
$+$ d ppl. », le tiret de fin de ligne après « Aqq » pouvant être un tiret de
remplissage, et la transcription le laisse tel quel ; mais la suite (vue 34)
écrit « $-$ Z sol. in A $+$ Dppl. », qui est $D - Zx$, et c'est ce que demande
l'équation de départ. Cette lecture suit la suite.} On ajoute aux deux membres
$B^4 - 2B^2x^2$, pour un $B$ quelconque :
\[ (x^2 - B^2)^2 = B^4 - 2B^2x^2 - Zx + D . \]
Soit $N^2 = 2B^2$ ; on égale les deux membres à $N^2y^2$. D'un côté
$x^2 - B^2 = Ny$ : une parabole. De l'autre
\[ \frac{B^4}{N^2} - x^2 - \frac{Z}{N^2}\,x + \frac{D}{N^2} = y^2 : \]
un cercle. Le calcul est juste ; c'est pour que le coefficient de $x^2$ soit
$-1$ après la division, et que la seconde courbe soit un cercle, que $N^2$ est
pris égal à $2B^2$.\footnote{Le feuillet écrit, au début de la vue 34, « æquale
Bqq $-$ Bq in Aq $-$ Z sol. in A $+$ Dppl. » : il faut « Bq in Aq bis »,
$2B^2x^2$, comme au membre de gauche et comme le demande la fraction qui suit.
À la fin de la vue 33, entre « Bqq, » et « Bq in Aq bis », aucun signe : il
faut un moins. La lettre N est écrite comme un n minuscule.}

La méthode s'étend à tous les cas : après la suppression du terme en $x^3$,
toute quartique s'écrit $x^4 = px^2 + qx + r$ ; toute cubique, débarrassée de
son terme en $x^2$ par la méthode de Viète puis multipliée par $x$, devient une
quartique sans terme en $x^3$. « Il faut seulement veiller, dans la seconde
équation, à ce que $x^2$ d'un côté et $y^2$ de l'autre aient des signes
contraires. » L'autre exemple du feuillet est $x^4 = px^2 - qd$, où $p$ est le
plan que le feuillet écrit « Zpl. » et $q$ le solide qu'il écrit « Z Sol. », et
$d$ une droite.\footnote{Le feuillet appelle $Z$ les deux grandeurs, l'une plane
et l'autre solide. « Z Sol. in d » est lu avec la lettre d, et fait de $qd$ un
terme constant : l'équation est alors bicarrée. Si ce « d » était une $A$ mal
formée, le terme $-qx$ passerait dans l'équation du cercle comme dans l'exemple
précédent, et la construction serait la même ; cette lecture garde la leçon.}
Ajoutant $B^4 - 2B^2x^2$ :
\[ (x^2 - B^2)^2 = B^4 - (2B^2 - p)\,x^2 - qd , \]
et l'on pose $N^2 = 2B^2 - p$, « la différence entre $2B^2$ et le plan » : la
parabole $x^2 - B^2 = Ny$ et le cercle
$\dfrac{B^4}{N^2} - x^2 - \dfrac{qd}{N^2} = y^2$. Il faut que $2B^2$ surpasse
$p$ : sinon « $x^2$ ne serait pas affecté du signe moins, et au lieu d'un cercle
nous trouverions une hyperbole ; à quoi le remède est prompt, car nous prenons
$B$ à volonté ». Tout cela est exact.\footnote{Le feuillet ajoute : « il est
constant, par la méthode des lieux, qu'il naît toujours un cercle d'une
équation dans l'un des membres de laquelle un carré inconnu est affecté du signe
$+$, dans l'autre un autre carré inconnu du signe $-$ ». « Toujours » est trop
dire : il faut encore que les deux carrés aient le même coefficient — c'est le
cas ici, après la division par $N^2$ — et que l'angle soit droit, faute de quoi
c'est une ellipse ; l'Isagoge le disait elle-même (vues 19 et 21). Et le cercle
peut être vide. Sur le feuillet, « plus » et « minus » sont barrés et remplacés
par les signes $+$ et $-$.}

Appliquée aux deux moyennes proportionnelles : $x^3 = B^2d$, donc
$x^4 = B^2dx$ ; ajoutant $B^4 - 2B^2x^2$,
\[ (x^2 - B^2)^2 = B^4 + B^2dx - 2B^2x^2 ; \]
avec $N^2 = 2B^2$, la parabole $x^2 - B^2 = Ny$ et le cercle
\[ \frac{B^2}{2} + \frac{d}{2}\,x - x^2 = y^2 , \]
ce que le feuillet écrit « Bq ½ $+$ d ½ in A $-$ Aq æquale Eq ». C'est juste.
Ici aussi la multiplication par $x$ a introduit une intersection parasite :
en éliminant $y$, on trouve $x\,(x^3 - B^2d) = 0$, et les deux courbes se
coupent au point $\bigl(0, -B/\sqrt2\bigr)$, qu'il faut écarter ; l'autre point
commun donne la plus grande moyenne.\footnote{« AC. » est l'abréviation de « A
cubus ». Les deux petits traits après « Aq » et « Bq » valent « bis ».}

\subsection*{La conclusion}

« Qui aura remarqué ces choses tentera en vain de tirer des lieux plans,
c'est-à-dire d'expédier par des droites et des cercles, la question des deux
moyennes, la trisection de l'angle et les semblables. » L'énoncé est vrai, mais
le feuillet ne le démontre pas, et ce qui précède ne le démontre pas non plus :
avoir résolu une question par des coniques ne prouve pas qu'on ne puisse la
résoudre par des cercles. La démonstration demande de savoir qu'un nombre
constructible à la règle et au compas est de degré une puissance de 2, et que
$x^3 = 2$, comme l'équation de la trisection d'un angle quelconque, est une
cubique irréductible ; elle a été donnée par Wantzel en 1837.\footnote{Ce
renseignement est général ; cette lecture ne prête à l'écrivain aucune
connaissance de ce résultat.}

\section*{Le lieu à trois droites (ff. 15r–16v, vues 37 à 40)}

\pagerange{37}{40}
En tête de la vue 37, d'une encre plus pâle et d'un trait plus fin que le
texte, peut-être d'une autre main : « Mr. fermat ». L'inscription rapporte la
pièce à Fermat ; elle ne dit pas de qui est l'écriture. Le texte s'adresse à
quelqu'un (« Superflua præsertim tibi non addimus », nous ne t'ajoutons pas le
superflu), et, au dos, le long du bord, une adresse est écrite en travers :
« Pour Mons.r Carcaui. rue michel le Conte au milieu », le nom lu avec
doute.\footnote{Le feuillet a été plié et porté comme une lettre. L'adresse ne
nomme pas l'expéditeur ; sa main paraît plus libre que celle du texte, et la
transcription ne décide pas si c'est la même.} La figure, que la page de texte
n'a pas, est seule au verso du feuillet suivant (vue 40).

\subsection*{L'énoncé}

Trois droites données de position, $AM$, $MB$, $BA$, forment un triangle. D'un
point $O$ on mène à ces droites $OE$, $OI$, $OD$ (vers $AM$, $MB$ et $AB$
respectivement) sous des angles donnés. Si le rapport du rectangle
$EO\cdot OD$ au carré de $OI$ est donné, le point $O$ est sur une section
conique. C'est le lieu à trois droites, ou problème de Pappus pour trois
droites.

La forme moderne. Soient $\ell_1$, $\ell_2$, $\ell_3$ des fonctions affines qui
s'annulent sur $AM$, $AB$ et $MB$. Une droite menée à une droite fixe sous un
angle fixe est proportionnelle à la distance du point à cette droite ; la
condition s'écrit donc $\ell_1\ell_2 = c\,\ell_3^2$, $c$ constant, et c'est une
équation du second degré : une conique. Toutes les coniques de cette famille
passent par $M$ et par $B$, où elles touchent respectivement $AM$ et $AB$ ; le
lieu cherché est celle d'entre elles qui passe par un point donné.

\subsection*{La construction du feuillet}

Couper $MB$ en deux en $q$ et joindre $Aq$. Par $O$ mener $foc$ parallèle à
$MB$ ($f$ sur $AM$, $c$ sur $AB$) et $On$ parallèle à $MA$ ($n$ sur $MB$). Les
triangles $Oef$, $Odc$, $Oin$ ont des angles donnés : ils sont de forme donnée,
et les rapports $OE/Of$, $OD/Oc$, $OI/On$ sont donnés. Donc
$(fO\cdot Oc)/On^2$ est donné, et, comme $On = fM$ (parallélogramme
$fMnO$), le rapport
\[ \frac{fO\cdot Oc}{fM^2} \]
est donné. Soit $V$ le point de $Aq$ tel que, $VR$ étant menée parallèlement à
$MB$ jusqu'à $AM$, $VR^2/RM^2$ soit égal à ce rapport ; puis la conique de
diamètre $Aq$, passant par $V$, que $MA$ et $AB$ touchent en $M$ et en $B$.
Elle passe par $O$.\footnote{« Si \emph{fertur} Aq in V » : le premier mot se
lit « fertur » ou « sertur » ; le sens demande « secetur », que $Aq$ soit
coupée en $V$, et c'est ainsi que cette lecture l'entend. Le feuillet écrit
une fois « foe » pour « foc », et biffe un mot court après « propter datum ».
Les lettres de la figure sont écrites tantôt en capitales, tantôt en
minuscules ; la lettre $V$ est le petit u du copiste.} Un tel point $V$ existe
et il est unique : quand $V$ va de $q$ vers $A$, $VR^2/RM^2$ décroît
continûment de l'infini à $0$. Le feuillet ajoute que, selon la position de
$V$, la conique sera une parabole, une hyperbole ou une ellipse ; précisément,
c'est une parabole si $V$ est le milieu de $Aq$, une ellipse si $V$ est entre
ce milieu et $q$, une hyperbole s'il est entre ce milieu et $A$.\footnote{Ce
critère est un complément de cette lecture. Avec $A = (0,1)$, $M = (-1,0)$,
$B = (1,0)$ et $V = (0,v)$, la conique est
$X^2 + \bigl(\tfrac1{v^2} - \tfrac2v\bigr)Y^2 + 2Y - 1 = 0$, de discriminant
$4(2v-1)/v^2$.}

\subsection*{L'analyse du feuillet, et Apollonius III.16}

Le feuillet raisonne ainsi. Supposons que la conique passe par $O$ et coupe
encore la parallèle $foc$ en $p$. La tangente en $V$ est $VR$, parallèle aux
ordonnées du diamètre $Aq$. Par la seizième proposition du troisième livre
d'Apollonius,
\[ \frac{pf\cdot fO}{fM^2} = \frac{VR^2}{RM^2} ,\]
et ce rapport vaut, par construction, $(fO\cdot Oc)/fM^2$. Donc
$pf\cdot fO = fO\cdot Oc$, $pf = Oc$, et, en ôtant $Op$, $fO = pc$. « Ce qui est
bien le cas » : $Aq$ coupant $MB$ en deux, elle coupe aussi $fc$ en deux en
$x$, et, étant un diamètre, la corde $Op$ aussi ; $fx = xc$ et $Ox = xp$
donnent $fO = pc$.\footnote{« Ex sexta decima propositione 3i. Apoll. » :
« decima » est écrit au-dessus de « sexta », et « Ideoque est rectangulum » est
barré juste avant. La leçon finale est la proposition 16 du livre III. La
lettre $x$ est celle de la figure, non une coordonnée.}

La proposition invoquée est, dans l'énoncé qu'on lui donne d'ordinaire, que,
si deux tangentes à une conique se rencontrent, et qu'une parallèle à l'une
d'elles coupe la conique en deux points et l'autre tangente en un point, le
rectangle des segments compris entre ce point et la conique est au carré du
segment de tangente qu'elle retranche comme les carrés des deux
tangentes.\footnote{C'est l'énoncé courant de la proposition ; il n'a été
confronté à aucune édition d'Apollonius. Il est vérifié ci-dessous.} C'est
bien ce dont le feuillet se sert, avec les tangentes $RV$ et $RM$ issues de
$R$, la parallèle $fOp$ à $RV$, et la tangente $AM$ coupée en $f$. La
proposition est vraie, et se démontre ainsi. Soit $Q = 0$ l'équation de la
conique et $q$ sa partie quadratique. Sur la droite menée d'un point $P$ dans
une direction $u$, $Q(P + tu)$ est un trinôme en $t$ de coefficient dominant
$q(u)$ et de terme constant $Q(P)$ ; le produit des distances (signées) de $P$
aux deux points d'intersection est donc $Q(P)/q(u)$, et, pour une tangente, le
carré de la distance au point de contact. Si $u$ est la direction de $MB$ et
$w$ celle de $AM$,
\[ \frac{fO\cdot fp}{fM^2} = \frac{Q(f)/q(u)}{Q(f)/q(w)} = \frac{q(w)}{q(u)}
   = \frac{RV^2}{RM^2} : \]
le rapport ne dépend que des directions. Le même calcul redonne la
proposition en toute généralité.

\subsection*{La synthèse que le feuillet laisse à faire}

Le raisonnement du feuillet est une analyse : il suppose ce qu'il faut prouver
et en tire une chose vraie, ce qui ne prouve rien. Il le sait, et conclut :
« nec est difficilis ab analysi ad Synthesim per demonstrationem ducentem ad
impossibile regressus » — le retour de l'analyse à la synthèse, par une
démonstration par l'absurde, n'est pas difficile. Il ne la donne pas. Voici une
synthèse directe, sans absurde. Sur la droite $fc$, prenons le milieu $x$
pour origine ; $x$ étant sur le diamètre conjugué de la direction $u$, le
trinôme n'y a pas de terme du premier degré, et, pour le point d'abscisse $t$,
\[ \frac{Q(x + tu)}{q(u)} = \frac{Q(x)}{q(u)} + t^2 . \]
En $f$, par la tangente $AM$ et la proposition précédente, cette quantité vaut
$fM^2\cdot q(w)/q(u) = fM^2\cdot RV^2/RM^2 = fO\cdot Oc$ par construction. Donc
$Q(x)/q(u) = fO\cdot Oc - fx^2$, et en $O$ :
\[ \frac{Q(O)}{q(u)} = fO\cdot Oc - fx^2 + xO^2
   = fO\cdot Oc - (fx - xO)(fx + xO) = fO\cdot Oc - fO\cdot Oc = 0 , \]
puisque $fx - xO = fO$ et $fx + xO = xc + xO = Oc$. Le point $O$ est sur la
conique.\footnote{Les longueurs sont prises avec leur signe sur la droite
$fc$, ce qui rend le calcul valable quelle que soit la position de $O$ par
rapport à $x$. Le texte s'arrête sur « regressus. », au tiers de la vue 38.}

\section*{L'usage des racines secondes, envoyé à Carcavi (ff. 17r–18r, vues 41 à 43)}

\pagerange{41}{43}
Le titre, en quatre lignes : « Nouus Secundarum et Vlterioris ordinis Radicum
In Analyticis Vsus A Domino de Fermat ad Dominum De Carcaui Die 20a. Aprilis
Anno 1650 Missus ».\footnote{« Nouvel usage, en analyse, des racines secondes
et d'ordre ultérieur, envoyé par M. de Fermat à M. de Carcavi le 20 avril
1650. » Le premier mot, « Nouus », est lu avec doute : sa dernière lettre a la
forme d'un e, et « Nouus » est proposé d'après l'accord avec « Vsus » et
« Missus ». C'est la page, non le catalogue, qui nomme ici Fermat comme auteur,
Carcavi comme destinataire, et qui donne cette date.} Une « racine seconde »
est une seconde inconnue. Le programme : ramener les racines secondes et
ultérieures aux premières, ce qui est « du plus grand moment en algèbre », par
un seul fondement, l'« analogie d'une égalité double », à itérer autant qu'il
le faut.\footnote{« Redductio », avec deux d, est l'orthographe du feuillet.}

\subsection*{La méthode sur l'exemple}

Deux équations en $x$ (la racine première, $A$ du feuillet) et $y$ (la racine
seconde, $E$), où $Z$ est un solide et $N^2$ un plan\footnote{« Item B In A
$+$ E q $+$ D in E » : le « $+$ » après « E q » est un petit signe surélevé
qu'on pourrait prendre pour le t d'une abréviation ; la reprise plus bas,
« $BA + E^2 + DE$ », le confirme.} :
\[ x^3 + y^3 = Z, \qquad Bx + y^2 + Dy = N^2 . \]
Premier précepte : dans chaque équation, faire passer d'un même côté les termes qui
contiennent $y$. Posons $P = Z - x^3$ et $Q = N^2 - Bx$ :
\[ P = y^3, \qquad Q = y^2 + Dy . \]
Second précepte : écrire la proportion $P : y^3 = Q : (y^2 + Dy)$ — les deux
rapports valent $1$ — et égaler le produit des extrêmes à celui des moyens :
\[ P\,(y^2 + Dy) = Q\,y^3 , \]
soit $Zy^2 - x^3y^2 + ZDy - x^3Dy = N^2y^3 - Bxy^3$.\footnote{Le feuillet
écrit « $N^2E^2$ », l'exposant 2 bien net à côté du 3 de « $BAE^3$ » : le
produit des moyens donne $N^2E^3$, et la ligne suivante, divisée par $E$, a
bien $N^2E^2$. La transcription laisse la faute ; cette lecture la corrige. Le
signe devant $A^3DE$ est traversé d'un petit trait vertical ; le calcul demande
un moins.} Le second et le quatrième terme de la proportion contenant $y$, tout
se divise par $y$ :
\[ P\,(y + D) = Q\,y^2 . \]
« Cette nouvelle équation sera plus basse d'un degré au moins, quant à la
racine seconde, que la plus haute des deux proposées » : degré 2 en $y$ au lieu
de 3. On recommence avec elle, mise sous la forme $PD = Qy^2 - Py$, et avec
$Q = y^2 + Dy$ :
\[ PD\,(y^2 + Dy) = Q\,(Qy^2 - Py) ,
   \qquad\text{puis, divisé par } y, \qquad
   PD\,(y + D) = Q\,(Qy - P) . \]
Le feuillet développe ces deux lignes terme à terme, et tous ses termes sont
justes. Cette fois $y$ n'apparaît plus qu'au premier degré, « sous le côté », et
une seule division, l'« application », le donne :
\[ y = \frac{ZD^2 - x^3D^2 + N^2Z - N^2x^3 - BZx + Bx^4}
            {N^4 - N^2Bx - BxN^2 + B^2x^2 - ZD + x^3D}
     = \frac{P\,(Q + D^2)}{Q^2 - DP} . \]
Le feuillet écrit ce quotient tel quel, avec les deux termes $-N^2BA$ et
$-BAN^2$ séparés au dénominateur ; la factorisation est de cette lecture.

Dernier précepte : remplacer $y$ par cette valeur dans l'une des deux
équations, de préférence la plus basse, pour que le degré ne monte pas ;
« il est manifeste que la racine seconde a entièrement disparu, qu'on est
arrivé à une équation libre de toute irrationalité, et que la méthode est
générale ». S'il y a plus de deux inconnues, on ramène de même les troisièmes
aux premières et secondes, puis les secondes aux premières.

\subsection*{Ce que la méthode calcule}

Chaque « analogie » remplace les deux équations par une combinaison à
coefficients polynomiaux — le produit croisé est
$(y^2 + Dy)\cdot(P - y^3) - y^3\cdot(Q - y^2 - Dy)$ — qui est encore vérifiée
par toute solution commune et dont le degré en $y$ a baissé d'une unité. Au
bout, l'équation en $x$ seul est ce qu'on appelle aujourd'hui l'éliminant, et
la théorie de l'élimination lui donne une forme canonique, le résultant. Ici,
le résultant en $y$ des deux équations est
\[ \operatorname{Res}_y = P^2 + D^3P + 3DPQ - Q^3 , \]
de degré 6 en $x$, comme le veut le produit des degrés, $3\times 2$.
L'équation finale de la méthode, obtenue en portant la valeur de $y$ dans
$Q = y^2 + Dy$ et en chassant le dénominateur, est de degré 8 : à un signe
près, c'est $(N^2 - Bx)^2\cdot\operatorname{Res}_y$. Elle contient donc le
facteur parasite $(N^2 - Bx)^2$, dont la racine $x = N^2/B$ n'est pas, en
général, une solution du problème.\footnote{Ces calculs sont de cette lecture,
et ont été vérifiés par un calcul symbolique. Le nom de résultant est attaché
ici parce que l'éliminant de la méthode coïncide avec lui à ce facteur près ;
la procédure du feuillet n'est pas le déterminant de Sylvester, et la lecture
ne l'appelle pas ainsi.}

Trois précisions rendent la méthode exacte. Les divisions par $y$ perdent les
solutions où $y = 0$, qui n'existent ici que si $Z = x^3$ et $N^2 = Bx$ à la
fois. L'« application » suppose le dénominateur non nul : en une solution
commune, $Q^2 - DP = y^2(y^2 + Dy + D^2)$, qui ne s'annule pas pour $y$ réel
non nul. Enfin la substitution finale peut introduire des racines étrangères,
comme on vient de le voir. La division euclidienne, en $y$, du premier
polynôme par le second aurait donné l'équation du premier degré en une seule
étape, $y = (P + DQ)/(Q + D^2)$ ; les deux expressions de $y$ coïncident en
toute solution commune. La méthode est générale au sens où elle élimine
toujours ; pour qu'elle soit exacte, il faut en contrôler les facteurs perdus
et ajoutés, ce que fait le résultant.

\section*{L'« Appendix ad superiorem methodum » : radicaux, problèmes abondants, le cône (ff. 18v–20v, vues 44 à 48)}

\pagerange{44}{48}
\subsection*{Les radicaux}

« C'est à la méthode précédente qu'on doit l'expurgation parfaite et absolue
des irrationalités en algèbre » ; le remède de Viète, le seul connu jusqu'ici,
n'y suffit pas.\footnote{Les mots qui désignent ce remède, « Symetria
climatismus Vietæa », sont lus lettre à lettre, et « huic » avec doute ;
l'accord demande un nom féminin que la page ne donne pas clairement. Cette
lecture ne les interprète pas.} Proposons par exemple
\[ \sqrt[3]{Bx^2 - x^3} + \sqrt{x^2 + Zx} + \sqrt[4]{B^3x - x^4}
   + \sqrt{Gx - x^2} = N : \]
comment l'« analyste viétéen » s'en tirera-t-il ? « Tandis que croîtra le
travail, ne croîtra-t-il pas la difficulté ? »\footnote{La lettre qui ouvre la
dernière racine se lit G ou D. Le mot après « Extricabit » se lit « sese » ou
« secet ».}

L'exemple traité\footnote{La première racine est écrite d'abord « Lat. Cub.
(ZA$^2$) », puis, quatre lignes plus bas, « Lat. Cub. (ZA$^2$ $-$ A$^3$) », et
c'est cette seconde forme que suit tout le calcul ; cette lecture la suit.} :
\[ \sqrt[3]{Zx^2 - x^3} + \sqrt[3]{x^3 + B^2x} = D . \]
On isole un radical, $D - \sqrt[3]{x^3 + B^2x}
= \sqrt[3]{Zx^2 - x^3}$, puis on nomme racine seconde chacun des autres :
$y = \sqrt[3]{x^3 + B^2x}$. En cubant,
\[ D^3 - 3D^2y + 3Dy^2 - y^3 = Zx^2 - x^3, \qquad y^3 = x^3 + B^2x , \]
deux équations en $x$ et $y$, sans radical, auxquelles s'applique la méthode ;
le feuillet s'arrête là. Cette lecture a achevé l'élimination. Si $u + v = D$
avec $u^3 = a$ et $v^3 = b$, alors $D^3 = a + b + 3uvD$, d'où
$(D^3 - a - b)^3 = 27abD^3$ ; ici $a + b = Zx^2 + B^2x$, et l'équation sans
radical est
\[ \bigl(D^3 - Zx^2 - B^2x\bigr)^3 = 27D^3\,(x^3 + B^2x)(Zx^2 - x^3) , \]
de degré 6 en $x$ (coefficient dominant $27D^3 - Z^3$) ; c'est bien le
résultant en $y$ des deux équations précédentes. Pour des racines cubiques
réelles, elle équivaut à l'équation de départ, sauf dans le cas où les deux
radicaux valent tous deux $-D$.\footnote{Car
$u^3 + v^3 - D^3 + 3uvD = (u + v - D)\,(u^2 + v^2 + D^2 - uv + uD + vD)$, et le
second facteur, égal à $\tfrac12\bigl[(u-v)^2 + (u+D)^2 + (v+D)^2\bigr]$, ne
s'annule que pour $u = v = -D$.}

La règle générale : s'il ne suffit pas de racines secondes, on nomme les
radicaux racines troisièmes, quatrièmes, et ainsi de suite, et l'on élimine la
dernière comme une seconde, en tenant les autres pour connues, puis la
précédente, jusqu'aux premières. « Il n'est donc aucune irrationalité que cette
méthode ne force à disparaître. »\footnote{« Emanuerit » et « exulerre » sont
lus lettre à lettre ; le sens attend « euanuerit » (qu'elle ait disparu) et
une forme d'« exulare » (partir en exil).} C'est exact : une équation où
figurent des radicaux de polynômes se ramène, en nommant chaque radical une
nouvelle inconnue et en éliminant, à une équation polynomiale ; celle-ci
contient toutes les solutions de la première, et peut en avoir d'autres,
venues des autres déterminations des racines.

Le feuillet ajoute que, les irrationalités disparues, l'artifice de Viète
pour les questions numériques donne des solutions aussi approchées qu'on veut,
« alors que, tant que durent les irrationalités, on ne saurait en aucune façon
obtenir de solutions approchées ». La première affirmation est juste ; la
seconde est un jugement de l'écrivain, que les méthodes numériques d'aujourd'hui
démentent : elles s'appliquent directement aux équations à radicaux.\footnote{La
construction de cette phrase, qui reprend « Imò » à la fin de la vue 45, est
donnée par la transcription telle qu'elle se lit ; une parenthèse fermante y
suit « Asymmetriæ » sans parenthèse ouvrante.}

\subsection*{Problèmes déterminés, locaux, « à la surface », abondants}

« Il s'est offert aussi à celui qui cherchait plus loin » une méthode pour la
connaissance des lieux à la surface, et pour les problèmes où il est donné plus
qu'il n'en faut.\footnote{« Semira » est écrit en un mot ; le sens le coupe
« se mira ».} Un problème à une seule inconnue est « déterminé » : on cherche
un point. Un problème à deux inconnues, qu'on ne peut réduire à une, est
« local » : on cherche une ligne. À trois inconnues, on cherche une surface :
ce sont les « lieux à la surface ». Aux premiers les données suffisent ; aux
seconds il en manque une, aux troisièmes deux. Et il se peut, à l'inverse,
qu'il y ait trop de données : le problème est alors « abondant ». Les
problèmes locaux, par opposition, l'écrivain les appelle « déficients ».
Aujourd'hui : un système de $m$ équations indépendantes en $n$ inconnues
définit, en général, un ensemble de dimension $n - m$ ; un problème abondant
est un système surdéterminé, $m > n$.

L'exemple : sur une droite donnée $AC$, un point $B$ tel que le rectangle
$AB\cdot BC$ soit donné et que la différence des carrés $AB^2 - BC^2$ soit donnée
aussi. Il y a une donnée de trop.\footnote{Le feuillet écrit ensuite « Ideoque
Solutio Ipsius quæstionis Exposuit », dont le dernier mot est lu tel quel et
n'est pas compris.} L'avantage est que la solution ne demande plus
d'extraction de racine : comme $AB + BC = AC$ et
$AB^2 - BC^2 = (AB - BC)(AB + BC) = s$, on a $AB - BC = s/AC$, et
\[ AB = \frac{AC^2 + s}{2\,AC} , \]
par une simple division ; la donnée du rectangle, $r$, n'est plus qu'une
condition de compatibilité, $r = \dfrac{AC^4 - s^2}{4\,AC^2}$.\footnote{Ce calcul
est de cette lecture ; le feuillet ne donne que l'énoncé.} Or ces problèmes
sont très fréquents, « surtout en physique et chez les artisans », et la
méthode les expédie tous « par une simple application, sans extraction de
racines, à quelque puissance que montent les équations ».

L'exemple en lettres :
\[ x^3 + B^2x = Z^2D, \qquad Gx - x^4 = B^2N , \]
où $G$ est un solide et $N$ un plan ; le problème est supposé abondant, les
deux équations ayant une racine commune. On traite la racine $x$ comme on
traitait la racine seconde, jusqu'à ce qu'elle soit donnée par une simple
division. Cette lecture a fait le calcul par divisions successives : en
remplaçant $x^4$ par $Z^2Dx - B^2x^2$, la seconde équation devient
$B^2x^2 + (G - Z^2D)\,x - B^2N = 0$, du second degré ; la première, divisée par
celle-ci, laisse une équation du premier degré, qui donne
\[ x = \frac{B^2\,\bigl(B^2DZ^2 - DNZ^2 + GN\bigr)}
            {B^6 + B^4N + (G - DZ^2)^2} . \]
C'est exact sous l'hypothèse que les deux équations ont une racine commune et
une seule : la racine commune est alors racine de leur plus grand commun
diviseur, qui est du premier degré, et s'exprime rationnellement par les
données. Si elles en avaient deux, le plus grand commun diviseur serait du
second degré, et la dernière division donnerait $0 = 0$.

\subsection*{Le cône coupé suivant un cercle}

« En guise de couronnement », le « fameux problème » : une ellipse et un point
hors de son plan étant donnés, couper par un plan la surface conique de sommet
ce point et de base cette ellipse, de sorte que la section soit un
cercle.\footnote{« Datũ », qualifiant le sommet, est lu avec doute.} Les
géomètres, dit le texte, prennent cinq points de l'ellipse, les joignent au
sommet, et cherchent un cercle qui rencontre ces cinq droites ; mais, les
points de l'ellipse étant en nombre infini, si l'on en prend six, le problème
devient abondant, et il en naît une égalité double qui donne l'inconnue par une
simple division. Le feuillet s'en tient à cette indication, sans calcul, et
cette lecture ne reconstitue pas la solution. Le problème a toujours des
solutions : le cône sur une ellipse est un cône du second degré, et ses
sections circulaires sont les sections par les plans parallèles à deux plans
fixes, ses plans cycliques, qui se confondent quand le cône est de
révolution.\footnote{Ce fait est général et n'est pas sur le feuillet. Si le
cône s'écrit $x^2/a^2 + y^2/b^2 = z^2/c^2$ avec $a > b$, les plans cycliques
sont les deux plans d'équation
$\bigl(\tfrac1{b^2} - \tfrac1{a^2}\bigr)y^2 = \bigl(\tfrac1{c^2} + \tfrac1{a^2}\bigr)z^2$.}

Le texte annonce enfin qu'on trouvera de même les diamètres et les axes d'une
courbe plane ou d'une surface donnée, de degré quelconque ; qu'on coupera une
surface conique dont la base est une parabole ou une ellipse « cubique »,
« quadrato-quadratique » ou de degré supérieur, de sorte qu'y paraisse telle
courbe qu'on voudra ; et « nous n'ajoutons rien sur les tangentes des courbes
et plusieurs autres usages de ce genre ». Ce sont des annonces, rapportées
comme telles.\footnote{Les deux mots brefs qui ouvrent la phrase sur la
surface conique, après une initiale en volute, se lisent lettre à lettre
« Oti di » ou « Vti dr » ; la construction attend un verbe comme « poterit ».
Ils restent illisibles. « parabole » est lu avec doute pour sa finale. La pièce
finit sur un trait de plume, sans signature.}

\section*{La lettre à Carcavi (ff. 21r–22r, vues 49 à 51)}

\pagerange{49}{51}
« Lettre de Mons.r De fermes / A Monsieur De Carcaui De Castres / Le 20e auril
1650. », de la main des vues 41 à 48, en français.\footnote{« De Castres » est
écrit à la suite du nom, séparé par un blanc : la page ne dit pas si c'est le
lieu d'où la lettre part ou une qualité du destinataire, et cette lecture ne
le dit pas davantage.} Elle commence : « Monsieur Je n'ay point voulu differer
a vous Enuoyer ma methode generale pour le desbrouilement des Asymmetries ».
L'épistolier envoie « mon original par pure paresse », demande qu'on le lui
renvoie « ou bien vne coppie », et ne réclame pour ses intérêts que « la
satisfaction … d'auoir deuoilé vne matiere qui n'estoit pas cognüe ». Il parle
à la fin de « mon escrit Latin », que ces éclaircissements doivent rendre moins
concis.\footnote{Que la lettre accompagne la pièce latine des vues 41 à 48 est
une inférence, tirée de sa date, de son destinataire et de sa matière ; la
transcription la fait avec la même réserve. L'épistolier écrit aussi : « En
voila trop pour vne seconde Lettre ». La lettre finit sur « \&c. », sans
formule finale ni signature.}

\subsection*{La tangente à une courbe irrationnelle}

« Si vous voulés auoir le plaisir tout entier », proposez de trouver la
tangente d'une courbe dont la propriété soit, en prenant cette fois $A$ pour
l'ordonnée et $E$ pour l'abscisse — ici $y$ et $x$\footnote{Le feuillet écrit
dans la deuxième racine « B pp$^{\text{l}}$ », B plano-plan, une grandeur de
dimension quatre nommée d'après $B$, que cette lecture écrit $B^4$ ; le terme
suivant, une lettre bouclée (D ?) avec un 2 repassé, puis « BA », est lu avec
doute, et son signe se confond avec un soulignement. Le feuillet dit que ces
quatre termes « sont, dans ce cas, des droites » : chaque radical est de
dimension un.} :
\[ \sqrt[3]{Z^2y - y^3} + \sqrt[4]{B^4 - D^2By + y^4} + \sqrt{By - y^2}
   + \sqrt[5]{y^5 - B^4y} = B + y - x . \]
Que fera là, demande-t-il, « La methode de Mons.r Des Cartes que vous sçauës estre
Infiniment plus Embarassée que la mienne si les Asymmetries ne sont ostées » ?
Sa réponse : en nommant chaque ligne irrationnelle racine seconde, tierce,
quarte, on arrive toujours à des égalités doubles, que l'on réitère jusqu'à ce
que la division simple ôte la dernière racine, puis l'avant-dernière ; il reste
une équation sans irrationalité entre $x$ et $y$, qui « representera la
proprieté specifique de la Courbe ». Alors « ma methode de Tangentibus » donne
la tangente « par La seule application », c'est-à-dire par une division, « ce
qui semble merueilleux ». Il n'ajoute pas l'opération entière, à cause de sa
longueur.\footnote{Que la « methode de Tangentibus » de la lettre soit celle des
vues 29 et 30, les feuillets ne le disent pas, et cette lecture ne le dit pas
non plus. Le jugement sur la méthode de Descartes est celui de l'épistolier ;
cette lecture ne l'examine pas.}

Ce qu'il affirme est exact. Pour une courbe $F(x,y) = 0$ polynomiale, la
tangente en un point régulier a pour pente $-F_x/F_y$, fraction rationnelle des
coordonnées du point : on la trouve sans extraire de racine. On remarquera
aujourd'hui qu'ici l'élimination n'est même pas nécessaire : l'abscisse $x$
figure au premier degré, la courbe est le graphe
$x = B + y - \sum(\text{radicaux})$, et la dérivée de chaque radical,
$\tfrac{d}{dy}\sqrt[k]{p(y)} = p'(y)/\bigl(k\,\sqrt[k]{p(y)}^{\,k-1}\bigr)$,
s'exprime rationnellement par $y$ et par la valeur du radical au point.

\subsection*{Le problème de Descartes}

La lettre poursuit : ce que ces méthodes donnent était inconnu « Jusqu'a
present, puisque Monsieur Des Cartes que Je nomme auec tout le respect qui est
deub a la memoire d'vn si grand et si merueilleux homme, proposoit comme vne
difficulté Insurmontable La question suiuante » : étant donnés quatre points,
et une courbe telle que, pour chacun de ses points, les droites menées aux
quatre points donnés fassent une somme donnée, trouver la tangente en un point
quelconque de cette courbe. « Ainsy que Je puis faire voir par vne de ses
Lettres » ; ce renvoi n'est pas vérifié ici.\footnote{La lettre parle de
Descartes comme d'un mort (« a la memoire ») ; c'est un renseignement général,
qui ne sert ici à rien dater, que Descartes est mort en février 1650. Un
premier « homme » est biffé ; une parenthèse fermante suit « deub » sans
parenthèse ouvrante.} L'épistolier affirme que ses méthodes jointes en donnent
la solution simple, et que « l'opération se fait en se jouant ».

La courbe s'appelle aujourd'hui une ellipse à quatre foyers (une
$n$-ellipse pour $n = 4$). Sa tangente se trouve en effet simplement : si
$P_1,\dots,P_4$ sont les points donnés et $P$ un point de la courbe, le gradient
de $\sum_i |P - P_i|$ est $\sum_i u_i$, où $u_i$ est le vecteur unitaire dirigé
de $P_i$ vers $P$ ; la normale en $P$ est dirigée par la somme de ces quatre
vecteurs unitaires, et la tangente lui est perpendiculaire — là où cette somme
n'est pas nulle et où $P$ n'est aucun des $P_i$. La construction se fait à la
règle et au compas. La lettre, elle, n'en donne pas l'opération.

Elle conclut : le principal effet de la nouvelle méthode « paroist aux
tangentes de toutte sorte de lignes a l'Infiny », qui s'y trouvent toujours
« par application simple », et ensuite « aux questions que J'appelle
abondantes, qui se resoluent aussy par la seule diuision et sans aucunes
Extractions de Racines ».

\end{document}
